% =============================================================================
% Apostila-guia — Git, GitLab e IA aplicada à Engenharia de Software
% Compilar com:  latexmk -xelatex apostila.tex
% =============================================================================
\documentclass[11pt,a4paper,oneside,openany]{book}
% =============================================================================
% marinha-comum.tex — identidade visual compartilhada pela apresentação e pela
% apostila da capacitação. Segue o Manual de Identidade Visual da Marinha do
% Brasil (dez/2025): cores institucionais (p. 23) e tipografia (pp. 24-30).
%
% Compilar sempre com XeLaTeX (fontspec). Os caminhos são relativos à pasta
% do documento (apresentacao/ ou apostila/).
% =============================================================================

% ---------- Caminhos ---------------------------------------------------------
\newcommand{\raizcapacitacao}{..}
\newcommand{\imagens}{\raizcapacitacao/assets/imagens}
\newcommand{\fontes}{\raizcapacitacao/assets/fontes}

% ---------- Cores institucionais (Manual, p. 23) -----------------------------
\usepackage{xcolor}
\usepackage{graphicx}
\definecolor{azulMB}{HTML}{050F41}   % Azul Marinha  · Pantone 2768C
\definecolor{ouroMB}{HTML}{FAB932}   % Amarelo Ouro  · Pantone 1235C
\definecolor{verdeMB}{HTML}{079551}  % Verde Brasil  · Pantone 355C
% Tons de apoio derivados das institucionais (fundos de caixas e diagramas)
\colorlet{azulClaro}{azulMB!8}
\colorlet{azulMedio}{azulMB!55}
\colorlet{ouroClaro}{ouroMB!18}
\colorlet{verdeClaro}{verdeMB!12}
\definecolor{cinzaTexto}{HTML}{3A3F55}
\definecolor{cinzaSuave}{HTML}{F3F4F8}
\definecolor{vermelhoAlerta}{HTML}{B3261E}
\colorlet{vermelhoClaro}{vermelhoAlerta!10}

% ---------- Tipografia (Manual, pp. 24-30) -----------------------------------
% Principal: Octin College (títulos) e alternativa Bebas Neue Pro são fontes
% comerciais. Usamos a Bebas Neue livre (OFL), da mesma família de desenho,
% nos títulos, e a Montserrat (tipografia de apoio oficial, OFL) no texto.
\usepackage{fontspec}
\setmainfont{Montserrat}[
  Path=\fontes/montserrat/, Extension=.otf,
  UprightFont=*-Regular, BoldFont=*-SemiBold,
  ItalicFont=*-Italic, BoldItalicFont=*-SemiBoldItalic]
\setsansfont{Montserrat}[
  Path=\fontes/montserrat/, Extension=.otf,
  UprightFont=*-Regular, BoldFont=*-SemiBold,
  ItalicFont=*-Italic, BoldItalicFont=*-SemiBoldItalic]
\setmonofont{FiraMono}[
  Path=\fontes/fira-mono/, Extension=.otf,
  UprightFont=*-Regular, BoldFont=*-Bold, Scale=0.9]
\newfontfamily\titulofonte{BebasNeue-Regular}[
  Path=\fontes/bebas-neue/, Extension=.ttf, LetterSpace=2.0]
\newfontfamily\montblack{Montserrat-Black}[Path=\fontes/montserrat/, Extension=.otf]
\newfontfamily\montlight{Montserrat-Light}[Path=\fontes/montserrat/, Extension=.otf]

% ---------- Pacotes comuns ----------------------------------------------------
\usepackage{tikz}
\usetikzlibrary{positioning,arrows.meta,calc,shapes.geometric,fit,backgrounds,shadows.blur,decorations.pathmorphing}
\usepackage{booktabs,tabularx,array}
\usepackage{fontawesome5}
\usepackage{listings}
\usepackage{csquotes}

% ---------- Código (listings) -------------------------------------------------
\lstdefinestyle{terminal}{
  basicstyle=\ttfamily\small\color{cinzaTexto},
  columns=fixed, basewidth=0.6em, keepspaces=true,
  breaklines=true, breakatwhitespace=true,
  showstringspaces=false, upquote=true,
  commentstyle=\color{verdeMB},
  keywordstyle=\color{azulMB}\bfseries,
  morekeywords={git,glab,claude},
  morecomment=[l]{\#},
  literate=
    {á}{{\'a}}1 {à}{{\`a}}1 {ã}{{\~a}}1 {â}{{\^a}}1
    {é}{{\'e}}1 {ê}{{\^e}}1 {í}{{\'i}}1 {ó}{{\'o}}1
    {õ}{{\~o}}1 {ô}{{\^o}}1 {ú}{{\'u}}1 {ç}{{\c c}}1
    {Á}{{\'A}}1 {É}{{\'E}}1 {Í}{{\'I}}1 {Ó}{{\'O}}1
    {Ú}{{\'U}}1 {Ç}{{\c C}}1 {Ã}{{\~A}}1 {Õ}{{\~O}}1
    {→}{{$\rightarrow$}}1 {—}{{---}}1 {·}{{$\cdot$}}1
}
\lstdefinestyle{saida}{style=terminal, keywordstyle=, commentstyle=,
  basicstyle=\ttfamily\small\color{cinzaTexto}}

% ---------- Diagramas de histórico Git (TikZ) ---------------------------------
\tikzset{
  commit/.style={circle, draw=#1, fill=white, line width=1.6pt,
    minimum size=7mm, inner sep=0pt, font=\small\bfseries\color{#1}},
  commit/.default=azulMB,
  commitm/.style={circle, draw=ouroMB!80!black, fill=ouroMB, line width=1.6pt,
    minimum size=8mm, inner sep=0pt, font=\small\bfseries\color{azulMB}},
  hist/.style={line width=1.8pt, color=#1},
  hist/.default=azulMB,
  rotulo/.style={rounded corners=2pt, fill=#1, text=white,
    font=\scriptsize\bfseries, inner sep=3pt},
  rotulo/.default=azulMB,
  caixa/.style={rounded corners=3pt, draw=#1, line width=1pt, fill=#1!8,
    align=center, font=\small, inner sep=5pt, text=azulMB},
  caixa/.default=azulMB,
  seta/.style={-{Stealth[length=3mm]}, line width=1.2pt, color=azulMB},
  setafina/.style={-{Stealth[length=2mm]}, line width=0.8pt, color=azulMedio},
}

% ---------- Marca do projeto (texto) -------------------------------------------
% Usada onde a versão interna tinha o logotipo da Marinha, que não entra na
% versão pública. \marca[cor das letras]{tamanho em pt}
\newcommand{\marca}[2][azulMB]{{\titulofonte\fontsize{#2}{#2}\selectfont
  \textcolor{#1}{GIT}\textcolor{ouroMB}{\,·\,}\textcolor{#1}{IA}\textcolor{ouroMB}{\,·\,}\textcolor{#1}{DOCS}}}

\usepackage{polyglossia}
\setdefaultlanguage[variant=brazilian]{portuguese}
\usepackage[top=28mm,bottom=25mm,left=25mm,right=22mm,headheight=14pt,headsep=10mm]{geometry}
\usepackage{microtype}
\usepackage{amssymb}
\usepackage{graphicx}
\usepackage{titlesec}
\usepackage{fancyhdr}
\usepackage{enumitem}
\usepackage{multicol}
\usepackage{pgfplots}
\pgfplotsset{compat=1.18}
\usepackage{dirtree}
\usepackage[most]{tcolorbox}
\usepackage[hidelinks,unicode]{hyperref}
\hypersetup{pdftitle={Git, GitLab e IA aplicada à Engenharia de Software — Apostila},
            pdfauthor={Ronivaldo D. Andrade}, pdfsubject={Capacitação interna — Marinha do Brasil}}
\urlstyle{same}

\color{cinzaTexto}
\setlength{\parindent}{0pt}
\setlength{\parskip}{6pt plus 1pt}
\linespread{1.12}
\setlist{itemsep=2pt, topsep=4pt}
\setlist[itemize,1]{label={\color{ouroMB}\small$\blacksquare$}}
\setlist[itemize,2]{label={\color{verdeMB}\scriptsize$\blacktriangleright$}}
\setlist[enumerate,1]{label={\color{azulMB}\bfseries\arabic*.}}

% ---------- Títulos -----------------------------------------------------------
\titleformat{\chapter}[display]
  {\color{azulMB}\raggedright\hyphenpenalty=10000}
  {\begin{tikzpicture}\node[inner sep=0pt, text=ouroMB, font=\titulofonte\fontsize{64}{64}\selectfont]{\thechapter};
   \fill[ouroMB] (0.05,-0.35) rectangle ++(2.6,0.12);\end{tikzpicture}}
  {2mm}
  {\titulofonte\fontsize{34}{36}\selectfont\MakeUppercase}
  [\vspace{1mm}]
\titleformat{name=\chapter,numberless}[display]
  {\color{azulMB}}{}{0pt}
  {\begin{tikzpicture}\fill[ouroMB] (0,0) rectangle ++(2.6,0.12);\end{tikzpicture}\\[3mm]\titulofonte\fontsize{34}{36}\selectfont\MakeUppercase}
\titlespacing*{\chapter}{0pt}{-6mm}{8mm}
\titleformat{\section}
  {\color{azulMB}\titulofonte\fontsize{21}{23}\selectfont}
  {\colorbox{azulMB}{\textcolor{white}{\,\thesection\,}}}{3mm}{\MakeUppercase}
\titlespacing*{\section}{0pt}{6mm}{2mm}
\titleformat{\subsection}
  {\color{azulMB}\large\bfseries}{\color{verdeMB}\thesubsection}{2.5mm}{}
\titlespacing*{\subsection}{0pt}{4mm}{1mm}
\setcounter{secnumdepth}{2}
\setcounter{tocdepth}{1}

% ---------- Cabeçalho e rodapé -------------------------------------------------
\pagestyle{fancy}
\fancyhf{}
\renewcommand{\headrulewidth}{0pt}
\renewcommand{\chaptermark}[1]{\markboth{#1}{}}
\fancyhead[L]{\footnotesize\color{azulMedio}\textsc{\leftmark}}
\fancyhead[R]{\marca{13}}
\fancyfoot[L]{\scriptsize\color{azulMedio}Git, GitLab e IA aplicada à Engenharia de Software \,·\, Apostila \,·\, Marinha do Brasil}
\fancyfoot[R]{\begin{tikzpicture}[baseline=(n.base)]\node[fill=azulMB, text=white, font=\footnotesize\bfseries, inner sep=3pt, minimum width=8mm](n){\thepage};\end{tikzpicture}}
\fancypagestyle{plain}{\fancyhf{}\fancyfoot[R]{\begin{tikzpicture}[baseline=(n.base)]\node[fill=azulMB, text=white, font=\footnotesize\bfseries, inner sep=3pt, minimum width=8mm](n){\thepage};\end{tikzpicture}}}

% ---------- Caixas (tcolorbox) ------------------------------------------------
\tcbuselibrary{listings,skins,breakable}
\tcbset{base/.style={enhanced, breakable, arc=2pt, boxrule=0pt, left=9pt, right=9pt, top=6pt, bottom=6pt,
  fonttitle=\bfseries\small, coltitle=white, before skip=8pt, after skip=10pt}}
\newtcolorbox{conceito}[1][]{base, colback=azulClaro, colframe=azulMB, borderline west={3pt}{0pt}{azulMB},
  title={\faBookOpen\enspace #1}, attach title to upper={\\[2pt]}, coltitle=azulMB}
\newtcolorbox{atencao}[1][Atenção]{base, colback=vermelhoClaro, colframe=vermelhoAlerta, borderline west={3pt}{0pt}{vermelhoAlerta},
  title={\faExclamationTriangle\enspace #1}, attach title to upper={\\[2pt]}, coltitle=vermelhoAlerta}
\newtcolorbox{dica}[1][Dica]{base, colback=verdeClaro, colframe=verdeMB, borderline west={3pt}{0pt}{verdeMB},
  title={\faLightbulb\enspace #1}, attach title to upper={\\[2pt]}, coltitle=verdeMB}
\newtcolorbox{analogia}[1][Analogia]{base, colback=ouroClaro, colframe=ouroMB, borderline west={3pt}{0pt}{ouroMB},
  title={\faComments\enspace #1}, attach title to upper={\\[2pt]}, coltitle=azulMB}
\newtcolorbox{exercicio}[2][]{base, colback=white, colframe=azulMB, boxrule=0.8pt, arc=3pt,
  title={\faDumbbell\enspace Exercício #2}, colbacktitle=azulMB, #1}
\newtcolorbox{verifique}{base, colback=cinzaSuave, colframe=cinzaSuave, title={\faCheckCircle\enspace Como saber se deu certo}, coltitle=verdeMB,
  attach title to upper={\\[2pt]}}

\tcbset{terminalbox/.style={base, colback=cinzaSuave, colframe=azulMB, boxrule=0pt, toprule=14pt,
  listing only, listing options={style=terminal}, left=6pt,
  overlay={\node[anchor=north west, text=white, font=\tiny\bfseries, inner sep=4pt] at (frame.north west) {\faTerminal\;TERMINAL #1};}},
  saidabox/.style={base, colback=white, colframe=azulMedio!50, boxrule=0.6pt, toprule=14pt,
  listing only, listing options={style=saida}, left=6pt,
  overlay={\node[anchor=north west, text=white, font=\tiny\bfseries, inner sep=4pt] at (frame.north west) {\faDesktop\;SAÍDA #1};}}}
% Sem argumento opcional: o "[" seria procurado já na linha do código e quebraria
% linhas que começam com "#". Para rótulo extra, use as versões com "t".
\newtcblisting{terminal}{terminalbox={}}
\newtcblisting{terminalt}[1]{terminalbox={#1}}
\newtcblisting{saida}{saidabox={}}
\newtcblisting{saidat}[1]{saidabox={#1}}
\lstset{style=terminal}

\newcommand{\cmd}[1]{\texttt{\color{azulMB}#1}}
\newcommand{\arq}[1]{\texttt{#1}}
\newcommand{\termo}[1]{\textbf{\color{azulMB}#1}}
\newcommand{\tecla}[1]{\tikz[baseline=(k.base)]\node[draw=azulMedio, rounded corners=2pt, inner sep=2pt, font=\footnotesize\ttfamily](k){#1};}

\begin{document}
\frontmatter
% ---------- Capa ---------------------------------------------------------------
\begin{titlepage}
\begin{tikzpicture}[remember picture, overlay]
  \fill[azulMB] (current page.south west) rectangle (current page.north east);
  \begin{scope}
    \clip (current page.north west) rectangle ([yshift=-0.42\paperheight]current page.north east);
    \node[anchor=north, inner sep=0pt] at (current page.north)
      {\includegraphics[height=0.42\paperheight]{\imagens/ilha_das_cobras.jpg}};
  \end{scope}
  \fill[ouroMB] ([yshift=-0.42\paperheight]current page.north west) rectangle ([yshift=-0.42\paperheight-1.6mm]current page.north east);
  \node[anchor=north west, text=ouroMB, font=\bfseries] (k) at ([xshift=22mm,yshift=-0.42\paperheight-16mm]current page.north west)
    {CAPACITAÇÃO INTERNA \,·\, APOSTILA-GUIA};
  \node[anchor=north west, text=white, align=left, inner xsep=0pt,
        font=\titulofonte\fontsize{44}{46}\selectfont] (t) at ([yshift=-3mm]k.south west)
    {GIT, GITLAB E IA APLICADA\\À ENGENHARIA DE SOFTWARE};
  \node[anchor=north west, text=white!82!azulMB, align=left, inner xsep=0pt, text width=15cm,
        font=\montlight\large] (s) at ([yshift=-4mm]t.south west)
    {Do primeiro commit à documentação de sistemas legados com um agente de IA: um guia para quem está começando};
  \node[anchor=south west, text=white, align=left, inner xsep=0pt, font=\normalsize] at ([xshift=22mm,yshift=22mm]current page.south west)
    {\textbf{Ronivaldo D. Andrade}\\[2pt]
     {\color{white!75!azulMB}\small Estagiário de Análise e Desenvolvimento de Sistemas}\\
     {\color{white!75!azulMB}\small Marinha do Brasil}\\[6pt]
     {\color{ouroMB}\small\bfseries Capacitação em 01/10/2026, das 10h às 11h}};
  \node[anchor=south east, inner sep=0pt] at ([xshift=-20mm,yshift=22mm]current page.south east)
    {\marca[white]{30}};
\end{tikzpicture}
\end{titlepage}

% ---------- Sobre esta apostila -------------------------------------------------
\chapter*{Sobre esta apostila}
\markboth{Sobre esta apostila}{}

Esta apostila acompanha a capacitação \textbf{\enquote{Git, GitLab e IA aplicada à Engenharia de Software}}, apresentada na Marinha do Brasil. A apresentação tem uma hora e é só expositiva. Aqui está tudo o que foi dito com mais calma, os detalhes que não cabem num slide e os \textbf{exercícios práticos}, para você aprender no seu ritmo.

Ela foi escrita para quem \textbf{nunca usou Git}. Se você já usa, pule direto para os capítulos \ref{cap:problema1} e \ref{cap:problema2}, que contam os dois problemas reais resolvidos no estágio.

\begin{center}
\renewcommand{\arraystretch}{1.35}
\begin{tabularx}{0.92\linewidth}{@{}>{\bfseries\color{azulMB}}l X@{}}
\toprule
Capacitação & Git, GitLab e IA aplicada à Engenharia de Software \\
Autor e instrutor & Ronivaldo D. Andrade \\
Função & Estagiário de Análise e Desenvolvimento de Sistemas \\
Organização & Marinha do Brasil \\
Data da capacitação & 01/10/2026, das 10h às 11h \\
Data de elaboração & 28/09/2026 \\
Material complementar & Apresentação de slides (\arq{apresentacao/apresentacao.pdf}) \\
\bottomrule
\end{tabularx}
\end{center}

\section*{Como ler esta apostila}

Ao longo do texto aparecem caixas coloridas. Cada cor tem um papel:

\begin{conceito}[Conceito]
Uma definição importante, para guardar.
\end{conceito}
\begin{analogia}
Uma comparação com o dia a dia, para a ideia ficar concreta.
\end{analogia}
\begin{dica}
Um atalho ou uma boa prática.
\end{dica}
\begin{atencao}
Um erro comum ou um cuidado que evita dor de cabeça.
\end{atencao}

\begin{terminal}
# Isto é um comando para digitar no terminal.
# As linhas que começam com # são comentários: não precisa digitar.
git status
\end{terminal}

\begin{saida}
Isto é o que o computador responde. O texto exato pode variar
conforme a versão e o idioma do seu Git.
\end{saida}

Quando um comando traz algo entre \arq{<sinais de menor e maior>}, como \arq{<URL\_DO\_REPOSITORIO>}, troque tudo, inclusive os sinais, pelo valor real.

\begin{atencao}[Antes de usar IA com código da instituição]
O capítulo \ref{cap:problema2} mostra como um agente de IA ajudou a documentar sistemas. Um agente desse tipo envia trechos do código para um serviço externo. \textbf{Só use com sistemas institucionais quando houver autorização e respeitando as normas de segurança da informação da sua OM.} Para praticar, use projetos pessoais ou de código aberto.
\end{atencao}

\tableofcontents
\mainmatter
\chapter{Antes de começar}
\label{cap:antes}

\section{Por que versionar?}

Quem já trabalhou num documento importante conhece esta pasta:

\begin{center}
\begin{tikzpicture}[doc/.style={draw=azulMedio, fill=cinzaSuave, rounded corners=2pt,
    font=\ttfamily\small, anchor=west, minimum width=9cm, inner sep=5pt}]
  \node[doc] at (0,0) {\faFile*[regular]\; relatorio.docx};
  \node[doc] at (0,-0.75) {\faFile*[regular]\; relatorio\_v2.docx};
  \node[doc] at (0,-1.5) {\faFile*[regular]\; relatorio\_v2\_final.docx};
  \node[doc] at (0,-2.25) {\faFile*[regular]\; relatorio\_v2\_final\_AGORA\_VAI.docx};
  \node[doc, draw=vermelhoAlerta, fill=vermelhoClaro] at (0,-3) {\faFile*[regular]\; relatorio\_v2\_final\_corrigido(1).docx};
\end{tikzpicture}
\end{center}

Ninguém sabe qual é a versão certa, quem mudou o quê, nem por quê. Agora imagine isso com um sistema de milhares de arquivos e várias pessoas mexendo ao mesmo tempo. É esse o problema que um \termo{sistema de controle de versão} resolve.

Com controle de versão, cada alteração fica registrada com:

\begin{itemize}
  \item \textbf{o que} mudou, linha a linha;
  \item \textbf{quem} mudou e \textbf{quando};
  \item \textbf{por que} mudou, numa mensagem escrita por quem fez;
  \item a possibilidade de \textbf{voltar} a qualquer estado anterior;
  \item várias pessoas trabalhando \textbf{em paralelo} sem apagar o trabalho umas das outras.
\end{itemize}

\begin{analogia}
Pense no Git como o \enquote{histórico de versões} de um documento compartilhado, só que muito mais poderoso: ele guarda o projeto inteiro, funciona sem internet e deixa cada pessoa trabalhar numa cópia própria antes de juntar tudo.
\end{analogia}

\section{O terminal, sem medo}

Quase tudo nesta apostila é feito no \termo{terminal} (também chamado de linha de comando, console ou \emph{shell}). É uma janela onde você digita uma ordem e aperta \tecla{Enter}.

\begin{itemize}
  \item \textbf{Linux}: procure por \enquote{Terminal} no menu de aplicativos.
  \item \textbf{Windows}: depois de instalar o Git, use o \textbf{Git Bash}, que vem junto. Ele aceita os mesmos comandos mostrados aqui.
\end{itemize}

Cinco comandos básicos ajudam a se localizar:

\begin{center}
\renewcommand{\arraystretch}{1.3}
\begin{tabularx}{0.9\linewidth}{@{}>{\ttfamily\color{azulMB}}l X@{}}
\toprule
pwd & mostra em que pasta você está \\
ls & lista os arquivos da pasta atual \\
cd pasta & entra numa pasta (\texttt{cd ..} volta uma pasta) \\
mkdir nome & cria uma pasta \\
cat arquivo & mostra o conteúdo de um arquivo \\
\bottomrule
\end{tabularx}
\end{center}

\section{Instalando o Git}

\begin{terminalt}{· LINUX (DEBIAN/UBUNTU)}
sudo apt update
sudo apt install git
\end{terminalt}

No \textbf{Windows}, baixe o instalador oficial em \url{https://git-scm.com/downloads} e aceite as opções padrão. Ao final, abra o \textbf{Git Bash}.

Para conferir se deu certo:

\begin{terminal}
git --version
\end{terminal}
\begin{saida}
git version 2.34.1
\end{saida}

Qualquer versão a partir da 2.28 funciona com todos os exemplos desta apostila. Se a sua for mais antiga, veja a dica da seção \ref{sec:primeiro-repo}.

\section{Dizendo ao Git quem é você}

Todo registro no Git leva o nome e o e-mail de quem o fez. Configure uma única vez por computador:

\begin{terminal}
git config --global user.name "Seu Nome"
git config --global user.email "seu.email@exemplo.com"
# faz a branch inicial de novos repositórios se chamar main
git config --global init.defaultBranch main
\end{terminal}

\begin{dica}
Use o mesmo e-mail cadastrado no GitLab da instituição. Assim o GitLab reconhece os seus commits e mostra a sua foto e o seu nome ao lado deles.
\end{dica}

Para ver o que está configurado:

\begin{terminal}
git config --list
\end{terminal}

\chapter{Git: o controle de versão}
\label{cap:git}

\begin{conceito}[Git]
O Git é um \textbf{sistema distribuído de controle de versão}. Ele registra a evolução de um projeto ao longo do tempo. \enquote{Distribuído} quer dizer que cada pessoa tem uma cópia completa do histórico no próprio computador: dá para trabalhar sem servidor e sem internet.
\end{conceito}

\begin{center}
\includegraphics[height=12mm]{\imagens/git-logo.pdf}
\end{center}

\section{Repositório e commit}

Um \termo{repositório} é uma pasta cujo histórico o Git está acompanhando. O Git guarda esse histórico numa subpasta escondida chamada \arq{.git}. Não mexa nela à mão.

Um \termo{commit} é uma \enquote{fotografia} do projeto num momento, acompanhada de quem a tirou, quando e por quê. Os commits se encadeiam e formam o \termo{histórico}:

\begin{center}
\begin{tikzpicture}
  \foreach \l/\x/\m in {A/0/Cria o projeto, B/2.6/Adiciona o login, C/5.2/Corrige o relatório, D/7.8/Atualiza o README}{
    \node[commit] (\l) at (\x,0) {\l};
    \node[font=\scriptsize, text width=2.3cm, align=center, below=2mm of \l] {\m};}
  \draw[hist] (A)--(B)--(C)--(D);
  \node[rotulo=verdeMB, above=4mm of D] (h) {main};
  \draw[setafina, verdeMB] (h) -- (D);
\end{tikzpicture}
\end{center}

Cada commit recebe um identificador único, um código como \arq{4203776}. Você vai ver esses códigos em todo lugar.

\section{As três áreas do Git}

Esta é a ideia que mais ajuda a entender os comandos. Um arquivo passa por três lugares:

\begin{center}
\begin{tikzpicture}[a/.style={caixa, minimum width=3.7cm, minimum height=2.2cm, text width=3.4cm}]
  \node[a] (w) at (0,0) {\textbf{Diretório de trabalho}\\[2pt]\scriptsize os arquivos como você os vê e edita};
  \node[a=ouroMB!85!black] (s) at (5,0) {\textbf{Área de preparação}\\[2pt]\scriptsize (\emph{staging}) o que vai entrar no próximo commit};
  \node[a=verdeMB] (r) at (10,0) {\textbf{Repositório}\\[2pt]\scriptsize o histórico de commits, guardado em \texttt{.git}};
  \draw[seta] (w.north east) to[bend left=25] node[above, font=\small\ttfamily] {git add} (s.north west);
  \draw[seta] (s.north east) to[bend left=25] node[above, font=\small\ttfamily] {git commit} (r.north west);
  \draw[seta, azulMedio] (r.south west) to[bend left=25] node[below, font=\small\ttfamily] {git restore / git switch} (w.south east |- r.south west);
\end{tikzpicture}
\end{center}

\begin{analogia}
É como preparar uma remessa. Você trabalha nos documentos na sua mesa (diretório de trabalho). Separa numa caixa os que vão ser enviados (área de preparação, com \cmd{git add}). Lacra a caixa e registra o envio no livro de protocolo (repositório, com \cmd{git commit}). O livro guarda cada remessa para sempre.
\end{analogia}

\section{Os primeiros comandos}
\label{sec:primeiro-repo}

\begin{terminal}
mkdir meu-projeto
cd meu-projeto
git init -b main          # transforma a pasta num repositório
\end{terminal}

\begin{dica}[Git anterior à versão 2.28]
Se o seu Git não aceitar \cmd{-b main}, use \cmd{git init} e, logo depois do primeiro commit, \cmd{git branch -m main}.
\end{dica}

Crie um arquivo e pergunte ao Git o que ele vê:

\begin{terminal}
echo "Olá, Git!" > saudacao.txt
git status
\end{terminal}
\begin{saida}
No ramo main

No commits yet

Arquivos não monitorados:
  (utilize "git add <arquivo>..." para incluir o que será submetido)
    saudacao.txt
\end{saida}

O Git viu o arquivo, mas ainda não o acompanha. Prepare e registre:

\begin{terminal}
git add saudacao.txt
git commit -m "Cria a saudação"
\end{terminal}
\begin{saida}
[main (root-commit) 4203776] Cria a saudação
 1 file changed, 1 insertion(+)
 create mode 100644 saudacao.txt
\end{saida}

\begin{dica}[Boas mensagens de commit]
Escreva o \textbf{porquê} da mudança, de um jeito que outra pessoa entenda daqui a um ano. \enquote{Corrige o cálculo do total quando não há itens} é útil. \enquote{ajustes} não é.
\end{dica}

Para consultar o histórico:

\begin{terminal}
git log --oneline         # uma linha por commit
git log --oneline --graph --all   # desenha as branches
git diff                  # o que mudou e ainda não foi preparado
\end{terminal}

\section{Branches}

\begin{conceito}[Branch]
Uma \textbf{branch} (ramo) é uma \textbf{linha de desenvolvimento} dentro do histórico. Ela permite trabalhar numa mudança de forma isolada, sem mexer na linha principal, a \arq{main}.
\end{conceito}

\begin{center}
\begin{tikzpicture}
  \node[commit] (a) at (0,0) {};  \node[commit] (b) at (1.5,0) {};
  \node[commit] (c) at (3,0) {}; \node[commit] (d) at (4.5,0) {}; \node[commit] (e) at (7.5,0) {};
  \draw[hist] (a)--(b)--(c)--(d)--(e);
  \node[commit=verdeMB] (f1) at (3,1.4) {}; \node[commit=verdeMB] (f2) at (4.5,1.4) {}; \node[commit=verdeMB] (f3) at (6,1.4) {};
  \draw[hist=verdeMB] (b) to[out=0,in=180] (f1); \draw[hist=verdeMB] (f1)--(f2)--(f3);
  \draw[hist=verdeMB, dashed] (f3) to[out=0,in=180] (e);
  \node[commit=ouroMB!85!black] (g1) at (4.5,-1.4) {}; \node[commit=ouroMB!85!black] (g2) at (6,-1.4) {};
  \draw[hist=ouroMB!85!black] (c) to[out=0,in=180] (g1); \draw[hist=ouroMB!85!black] (g1)--(g2);
  \node[rotulo, anchor=west] at ([xshift=2mm]e.east) {main};
  \node[rotulo=verdeMB, anchor=west] at ([xshift=2mm]f3.east) {feature/login};
  \node[rotulo=ouroMB!85!black, anchor=west] at ([xshift=2mm]g2.east) {bugfix/autenticacao};
\end{tikzpicture}
\end{center}

Branches servem para desenvolver funcionalidades isoladamente, corrigir bugs, experimentar, trabalhar em paralelo, evitar alterações diretas na \arq{main} e permitir revisão antes da integração. É comum usar prefixos que dizem o tipo do trabalho: \arq{feature/}, \arq{bugfix/}, \arq{hotfix/}.

\begin{atencao}[Branch não é cópia]
Branch \textbf{não} é uma cópia do projeto. É só um \textbf{ponteiro} que aponta para um commit e anda para a frente a cada novo commit. Por isso criar uma branch é instantâneo, mesmo num projeto enorme.
\end{atencao}

\begin{terminal}
git branch                        # lista as branches (* marca a atual)
git switch -c feature/despedida   # cria e já entra na nova branch
git switch main                   # volta para a main
\end{terminal}

\section{Merge: juntando linhas de trabalho}

O \termo{merge} traz as mudanças de uma branch para a branch em que você está.

\begin{terminal}
git switch main
git merge feature/despedida
\end{terminal}
\begin{saida}
Updating 4203776..afe16bb
Fast-forward
 saudacao.txt | 1 +
 1 file changed, 1 insertion(+)
\end{saida}

\enquote{Fast-forward} quer dizer que a \arq{main} não tinha nada novo, então o Git só avançou o ponteiro. Quando as duas branches andaram, o Git cria um \termo{commit de merge}, que tem dois \enquote{pais} e une as duas linhas.

\section{Conflitos}
\label{sec:conflitos}

\begin{conceito}[Conflito]
Um conflito acontece quando duas branches mudaram \textbf{a mesma região do mesmo arquivo} de formas diferentes, e o Git não tem como saber qual versão vale. Conflito \textbf{não é erro}: é o Git pedindo uma decisão humana.
\end{conceito}

Exemplo: na \arq{main}, a primeira linha virou \enquote{E aí, pessoal!}. Na \arq{feature/formal}, virou \enquote{Prezados, bom dia!}. Ao juntar:

\begin{terminal}
git merge feature/formal
\end{terminal}
\begin{saida}
Mesclagem automática de saudacao.txt
CONFLITO (conteúdo): conflito de mesclagem em saudacao.txt
Automatic merge failed; fix conflicts and then commit the result.
\end{saida}

O Git escreve as duas versões dentro do arquivo, separadas por marcadores:

\begin{saidat}{· CONTEÚDO DE saudacao.txt}
<<<<<<< HEAD
E aí, pessoal!
=======
Prezados, bom dia!
>>>>>>> feature/formal
Até logo!
\end{saidat}

\begin{center}
\renewcommand{\arraystretch}{1.3}
\begin{tabularx}{0.9\linewidth}{@{}>{\ttfamily\color{azulMB}}l X@{}}
\toprule
<<<<<<< HEAD & início da versão da branch em que você está \\
======= & separa as duas versões \\
>>>>>>> feature/formal & fim da versão que está chegando \\
\bottomrule
\end{tabularx}
\end{center}

Para resolver:

\begin{enumerate}
  \item abra o arquivo e deixe só o conteúdo que deve ficar (pode ser uma das versões ou uma combinação das duas);
  \item apague as três linhas de marcadores;
  \item marque como resolvido e registre:
\end{enumerate}

\begin{terminal}
git add saudacao.txt
git commit -m "Resolve o conflito da saudação"
\end{terminal}

\begin{dica}
Se você se perdeu no meio de um merge, \cmd{git merge --abort} desfaz tudo e volta ao estado de antes do merge. Nada se perde.
\end{dica}

\section{Trabalhando com um repositório remoto}

Até aqui tudo foi local. Para trabalhar em equipe, o repositório também fica num servidor, como o GitLab. Esse servidor é chamado de \termo{remoto}, e por convenção o principal se chama \arq{origin}.

\begin{center}
\renewcommand{\arraystretch}{1.3}
\begin{tabularx}{0.95\linewidth}{@{}>{\ttfamily\color{azulMB}}l X@{}}
\toprule
git clone <URL> & baixa um repositório remoto inteiro, já configurado \\
git remote -v & mostra os remotos configurados \\
git remote add origin <URL> & liga um repositório local a um remoto \\
git fetch & baixa as novidades do remoto \textbf{sem} mexer nos seus arquivos \\
git pull & baixa \textbf{e} já faz o merge na branch atual \\
git push origin <branch> & envia os seus commits para o remoto \\
\bottomrule
\end{tabularx}
\end{center}

\begin{analogia}[fetch × pull]
\cmd{git fetch} é olhar a caixa de correspondência e trazer as cartas para a mesa, sem abrir. \cmd{git pull} é trazer e já arquivar tudo na pasta. O \cmd{fetch} é mais seguro quando você quer ver antes o que chegou.
\end{analogia}

\section{Desfazendo coisas}

\begin{center}
\renewcommand{\arraystretch}{1.3}
\begin{tabularx}{0.95\linewidth}{@{}>{\ttfamily\color{azulMB}}l X@{}}
\toprule
git restore arquivo & descarta as mudanças \textbf{não preparadas} do arquivo \\
git restore --staged arquivo & tira o arquivo da área de preparação, mantendo as mudanças \\
git revert <commit> & cria um novo commit que desfaz outro (seguro em branch compartilhada) \\
git merge --abort & cancela um merge em andamento \\
\bottomrule
\end{tabularx}
\end{center}

\begin{atencao}
\cmd{git restore arquivo} apaga as mudanças que ainda não viraram commit. Não há como recuperá-las depois. Confira com \cmd{git diff} antes.
\end{atencao}

\section{Um aviso honesto}

\begin{center}
\begin{minipage}{0.3\linewidth}
\includegraphics[width=\linewidth]{\imagens/xkcd-1597-git.png}
\end{minipage}\hfill
\begin{minipage}{0.64\linewidth}
\small\itshape
\enquote{Isto é o Git. Ele controla trabalho colaborativo em projetos por meio de um belo modelo distribuído de árvore da teoria dos grafos.}\\[2pt]
\enquote{Legal. Como a gente usa?}\\[2pt]
\enquote{Não faço ideia. Só decora estes comandos de terminal e digita para sincronizar. Se der erro, salva teu trabalho em outro lugar, apaga o projeto e baixa uma cópia nova.}\\[6pt]
\upshape\footnotesize\color{azulMedio} xkcd \#1597, \enquote{Git}, de Randall Munroe (\url{https://xkcd.com/1597/}), licença CC BY-NC 2.5. Tradução livre.
\end{minipage}
\end{center}

A tirinha é engraçada porque é verdade para muita gente. O objetivo deste capítulo foi o contrário: entender o \textbf{modelo} (commits, ponteiros, três áreas), para que os comandos façam sentido e você nunca precise apagar o projeto.

\chapter{GitLab: colaboração em volta do Git}
\label{cap:gitlab}

\begin{conceito}[GitLab]
O GitLab é uma \textbf{plataforma de desenvolvimento e gerenciamento de software baseada no Git}. Ele hospeda repositórios e acrescenta, num só lugar, o que uma equipe precisa para colaborar: revisão de código, controle de permissões, Merge Requests, Issues, CI/CD, pipelines, releases, segurança e auditoria.
\end{conceito}

\section{Git, GitLab, GitHub e glab}

Esses nomes confundem porque são parecidos. Mas \textbf{Git e GitLab não são a mesma coisa}:

\begin{center}
\renewcommand{\arraystretch}{1.5}
\small
\begin{tabularx}{\linewidth}{@{}m{13mm} >{\bfseries\color{azulMB}}m{22mm} >{\raggedright\arraybackslash}X >{\raggedright\arraybackslash}X@{}}
\toprule
 & & \textbf{O que é} & \textbf{Principal finalidade} \\
\midrule
\includegraphics[height=5mm]{\imagens/git-logo.pdf} & Git & Sistema distribuído de controle de versão & Controlar o histórico do código \\
\includegraphics[height=7mm]{\imagens/gitlab-tanuki.pdf} & GitLab & Plataforma baseada em Git, com opção de instalação na própria infraestrutura & Hospedar, gerenciar e colaborar sobre projetos \\
\includegraphics[height=6mm]{\imagens/github-mark.pdf} & GitHub & Plataforma baseada em Git, oferecida principalmente como serviço em nuvem & Hospedar, gerenciar e colaborar sobre projetos \\
\centering\faTerminal & GitLab CLI (\texttt{glab}) & Ferramenta de linha de comando & Interagir com o GitLab pelo terminal \\
\bottomrule
\end{tabularx}
\end{center}

\begin{analogia}
O \textbf{Git} é o motor do carro. O \textbf{GitLab} e o \textbf{GitHub} são garagens com oficina, vistoria e controle de quem entra. O \textbf{glab} é o controle remoto do portão da garagem do GitLab.
\end{analogia}

O Git funciona sozinho, sem GitLab nem GitHub. Os comandos \cmd{git init}, \cmd{git add} e \cmd{git commit} do capítulo anterior não precisaram de servidor nenhum. As plataformas entram quando é preciso compartilhar e revisar.

A diferença entre GitLab e GitHub no dia a dia é pequena: os dois hospedam repositórios Git e permitem revisão. Os nomes mudam um pouco (o GitHub chama o Merge Request de \emph{Pull Request}). A diferença que mais importa para uma instituição é a próxima seção.

\section{GitLab Self-Managed}

O GitLab tem uma modalidade que a organização \textbf{instala e administra na própria infraestrutura}, chamada \termo{GitLab Self-Managed}. A instituição mantém a sua instância nos seus servidores ou em infraestrutura autorizada, sem mandar o código para fora.

\begin{center}
\begin{tikzpicture}
  \node[draw=azulMB, line width=1.2pt, rounded corners=6pt, minimum width=8.5cm, minimum height=4.4cm, fill=azulClaro] (om) at (0,0) {};
  \node[rotulo, anchor=north west] at ([xshift=2mm,yshift=-2mm]om.north west) {\faBuilding\; Infraestrutura da instituição};
  \node[caixa=verdeMB, minimum width=2.6cm] (srv) at (0,0.2) {\includegraphics[height=7mm]{\imagens/gitlab-tanuki.pdf}\\GitLab};
  \node[caixa, minimum width=1.8cm, font=\scriptsize] (u1) at (-2.8,-1.3) {\faUsers\\equipe};
  \node[caixa, minimum width=1.8cm, font=\scriptsize] (u2) at (2.8,-1.3) {\faServer\\backups};
  \draw[setafina,<->] (u1) -- (srv); \draw[setafina,<->] (u2) -- (srv);
  \draw[vermelhoAlerta, line width=1.2pt, dashed] (4.6,-2.2) -- (4.6,2.2);
  \node[font=\small\color{vermelhoAlerta}, align=center] at (6,0) {\faCloud\\internet\\externa};
\end{tikzpicture}
\end{center}

Isso dá à instituição mais controle sobre armazenamento dos dados, código-fonte, histórico, usuários, permissões, autenticação, auditoria, políticas de segurança e integração com outros sistemas.

\begin{atencao}[Dentro de casa não é sinônimo de seguro]
Não interprete Self-Managed como \enquote{se está dentro da instituição, está automaticamente seguro}. A segurança continua dependendo de configuração, controle de acesso, atualização do software, autenticação, política de senhas, backups, firewall, segmentação de rede, monitoramento, auditoria e gestão de vulnerabilidades. O ganho é outro: a instituição \textbf{mantém o ambiente sob a própria administração}.
\end{atencao}

\section{Branch protegida}

No GitLab, uma branch pode ser \termo{protegida}. Numa branch protegida, ninguém faz \cmd{git push} direto (ou só quem tem permissão), e a única forma de mudar o conteúdo é por \textbf{Merge Request}. É comum a \arq{main} ser protegida, para que nada chegue à versão principal sem revisão.

\section{Merge Request}

\begin{conceito}[Merge Request (MR)]
É uma \textbf{solicitação para incorporar as alterações de uma branch em outra}, normalmente acompanhada de revisão. No GitHub, o mesmo recurso se chama \emph{Pull Request}.
\end{conceito}

\begin{center}
\begin{tikzpicture}[passo/.style={caixa, minimum width=2.1cm, minimum height=1.25cm, font=\scriptsize}]
  \node[rotulo=verdeMB, font=\small\bfseries] (src) at (0,1.9) {feature/login};
  \node[rotulo, font=\small\bfseries] (dst) at (12.5,1.9) {main};
  \draw[seta, line width=2pt] (src) -- node[above, font=\small\bfseries\color{azulMB}] {Merge Request} (dst);
  \foreach \i/\ic/\t/\cor in {0/\faEye/Visualizar/azulMB, 1/\faCommentDots/Comentar/azulMB, 2/\faRedo/Pedir ajustes/azulMB, 3/\faVial/Testar/azulMB, 4/\faThumbsUp/Aprovar/azulMB, 5/\faCodeBranch/Merge/verdeMB}{
    \node[passo=\cor] (p\i) at (\i*2.5,0) {{\large\ic}\\[2pt]\t};}
  \foreach \i [evaluate=\i as \j using int(\i+1)] in {0,...,4}{\draw[setafina] (p\i) -- (p\j);}
\end{tikzpicture}
\end{center}

O fluxo é: a pessoa trabalha numa branch, envia a branch ao GitLab com \cmd{git push} e abre um Merge Request. A equipe então visualiza as alterações, revisa o código, comenta, pede ajustes, roda os testes automáticos (se houver pipeline), aprova e, por fim, faz o merge.

\begin{analogia}[Commit, branch e Merge Request]
\textbf{Commit} registra uma alteração. \textbf{Branch} cria uma linha de desenvolvimento. \textbf{Merge Request} pede que essa linha seja integrada, como um processo que só segue depois do \enquote{de acordo} de quem tem a competência para aprovar.
\end{analogia}

\subsection{Abrindo um Merge Request pela interface web}

Os nomes dos botões podem variar um pouco conforme a versão do GitLab instalada:

\begin{enumerate}
  \item envie a sua branch: \cmd{git push origin minha-branch};
  \item abra o projeto no GitLab. Logo após o push, ele costuma mostrar um aviso com o botão \textbf{Create merge request};
  \item se o aviso não aparecer, vá em \textbf{Code \faChevronRight\ Merge requests \faChevronRight\ New merge request} e escolha a branch de origem (a sua) e a de destino (em geral, a \arq{main});
  \item preencha um título claro e uma descrição do que mudou e por quê;
  \item indique revisores, se a equipe usar;
  \item clique em \textbf{Create merge request}.
\end{enumerate}

Depois disso, a aba \textbf{Changes} mostra linha a linha o que mudou, e a revisão acontece em comentários ali mesmo.

\section{O fluxo básico de trabalho em equipe}

\begin{center}
\begin{tikzpicture}[p/.style={caixa, minimum width=1.55cm, text width=1.35cm, minimum height=1.15cm, font=\scriptsize, inner sep=3pt}]
  \foreach \i/\t/\c in {0/{\ttfamily git\\pull}/azulMB, 1/Alterar\\código/azulMB, 2/{\ttfamily git\\status}/azulMB, 3/{\ttfamily git\\add}/azulMB,
                        4/{\ttfamily git\\commit}/azulMB, 5/{\ttfamily git\\push}/azulMB, 6/Merge\\Request/verdeMB, 7/Revisão\\e merge/verdeMB}{
    \node[p=\c] (n\i) at (\i*1.95,0) {\t};}
  \foreach \i [evaluate=\i as \j using int(\i+1)] in {0,...,6}{\draw[setafina] (n\i) -- (n\j);}
\end{tikzpicture}
\end{center}

\begin{terminal}
git switch main
git pull                              # começa sempre atualizado
git switch -c minha-branch            # trabalha numa branch própria
# ... realizar as alterações ...
git status
git add arquivo-alterado.py
git commit -m "Adiciona nova funcionalidade"
git push origin minha-branch
# depois: abrir o Merge Request no GitLab
\end{terminal}

\begin{dica}[Regra de ouro]
Comece sempre com \cmd{git pull} e mantenha a sua branch curta. Quanto mais tempo ela fica longe da \arq{main}, maior a chance de conflito.
\end{dica}

\begin{atencao}[\texttt{git add .}]
\cmd{git add .} prepara \textbf{tudo} o que mudou na pasta, inclusive arquivos que você não queria enviar. Confira com \cmd{git status} antes, ou adicione os arquivos pelo nome.
\end{atencao}

\chapter{GitLab CLI: o GitLab pelo terminal}
\label{cap:glab}

\begin{conceito}[GitLab CLI]
O GitLab CLI, usado pelo comando \cmd{glab}, é a ferramenta oficial para \textbf{interagir com o GitLab diretamente pelo terminal}: Merge Requests, Issues, projetos, pipelines, branches, releases e autenticação.
\end{conceito}

\begin{center}
\begin{tikzpicture}[b/.style={caixa, minimum width=3.6cm, minimum height=1.4cm}]
  \node[b] (g) at (0,0) {\textbf{Git}\\\scriptsize controle de versão};
  \node[b=verdeMB] (l) at (4.6,0) {\textbf{GitLab}\\\scriptsize plataforma de colaboração};
  \node[b=ouroMB!85!black] (c) at (9.2,0) {\textbf{glab}\\\scriptsize o GitLab no terminal};
\end{tikzpicture}
\end{center}

\section{Instalação}

As formas de instalar mudam conforme o sistema e a versão. Siga a página oficial: \url{https://gitlab.com/gitlab-org/cli}, seção \emph{Installation}. Para conferir:

\begin{terminal}
glab --version
\end{terminal}

\section{Autenticação}

O \cmd{glab} precisa saber em qual GitLab e com qual usuário trabalhar. Num GitLab instalado pela instituição, informe o endereço dele:

\begin{terminal}
glab auth login --hostname <endereco-do-gitlab>
\end{terminal}

O comando pergunta como você quer se autenticar. O caminho mais comum é um \termo{token de acesso pessoal}, gerado no GitLab em \textbf{Preferências do usuário \faChevronRight\ Access tokens}.

\begin{atencao}[Token é senha]
Um token de acesso dá acesso à sua conta. Não compartilhe, não envie por e-mail ou chat e não grave em arquivo de projeto. Dê a ele só as permissões necessárias e uma data de validade.
\end{atencao}

\section{Os comandos mais úteis}

\begin{center}
\renewcommand{\arraystretch}{1.3}
\begin{tabularx}{\linewidth}{@{}>{\ttfamily\color{azulMB}}l X@{}}
\toprule
glab mr create & cria um Merge Request a partir da branch atual \\
glab mr create -s origem -b destino & escolhe as branches de origem e de destino \\
glab mr list & lista os Merge Requests abertos do projeto \\
glab mr view & mostra o Merge Request da branch atual \\
glab mr diff & mostra as alterações do Merge Request \\
glab mr approve & aprova um Merge Request (se você tiver permissão) \\
glab mr merge & faz o merge (se você tiver permissão) \\
glab issue list & lista as Issues \\
glab ci status & mostra a situação do pipeline da branch atual \\
glab repo clone grupo/projeto & clona um projeto pelo nome \\
\bottomrule
\end{tabularx}
\end{center}

Exemplo completo de criação de Merge Request:

\begin{terminal}
glab mr create --source-branch desenvolvimento --target-branch main \
  --title "Integra o sistema ao repositório institucional" \
  --description "Une o histórico local ao histórico do GitLab."
\end{terminal}

Todo comando tem ajuda embutida: \cmd{glab mr create --help}.

\section{Terminal ou interface web?}

Tudo o que é versionamento e merge dá para fazer no Git, localmente. A parte de colaboração no GitLab dá para fazer com o \cmd{glab}. Mesmo assim, em ambiente institucional a \textbf{interface web continua importante} para revisão visual, aprovação, auditoria, acompanhamento dos pipelines e conversa da equipe.

\begin{atencao}[As regras valem no terminal também]
O \cmd{glab} não pula regras. Se a \arq{main} exige aprovação de um responsável, \cmd{glab mr merge} vai esperar essa aprovação como a interface web esperaria.
\end{atencao}

\chapter{Padrões de commit}
\label{cap:padroes}

No capítulo \ref{cap:git} vimos que um commit leva uma mensagem. Este capítulo trata de \textbf{como escrever essa mensagem} de um jeito que a equipe inteira (e as ferramentas) consigam ler.

\section{Por que padronizar}

Compare dois históricos do mesmo projeto:

\begin{saidat}{· SEM PADRÃO}
a41aded ajustes
c2524a8 arrumei
afe16bb agora vai
4203776 mudanças do joão
\end{saidat}
\begin{saidat}{· COM PADRÃO}
a41aded fix(login): corrige senha vazia
c2524a8 feat(relatorio): exporta em PDF
afe16bb docs: explica a instalação
4203776 test(login): cobre senha vazia
\end{saidat}

No segundo, dá para saber \textbf{o tipo} da mudança, \textbf{onde} ela aconteceu e \textbf{o que} foi feito, sem abrir um arquivo sequer. Um padrão de commit também permite automatizar tarefas: gerar notas de versão, decidir o próximo número de versão e barrar mensagens fora do combinado.

\begin{conceito}[A regra mais importante]
O melhor padrão é \textbf{o que a equipe combinou e segue}. Os padrões deste capítulo são os mais usados, mas consistência vale mais que qualquer escolha isolada.
\end{conceito}

\section{Conventional Commits}

O padrão mais difundido é o \termo{Conventional Commits} 1.0.0 (\url{https://www.conventionalcommits.org/pt-br/v1.0.0/}). A estrutura é:

\begin{terminalt}{· ESTRUTURA}
<tipo>[escopo opcional]: <descrição>

[corpo opcional]

[rodapé(s) opcional(is)]
\end{terminalt}

As regras principais da especificação:

\begin{itemize}
  \item o \textbf{tipo} é obrigatório e vem primeiro, seguido de \textbf{dois-pontos e espaço};
  \item o \textbf{escopo} é opcional: um substantivo entre parênteses que diz qual parte do sistema mudou, como \arq{fix(parser):};
  \item a \textbf{descrição} vem logo depois de \arq{: } e resume a mudança;
  \item o \textbf{corpo} é opcional e começa \textbf{uma linha em branco} depois da descrição, com quantos parágrafos forem necessários;
  \item os \textbf{rodapés} são opcionais e vêm uma linha em branco depois do corpo, no formato \arq{Token: valor} (por exemplo, \arq{Refs: \#123} ou \arq{Reviewed-by: Maria});
  \item uma \textbf{quebra de compatibilidade} (\emph{breaking change}) é indicada com \arq{!} antes dos dois-pontos, com um rodapé \arq{BREAKING CHANGE: <explicação>}, ou com os dois.
\end{itemize}

\begin{terminalt}{· EXEMPLO COMPLETO}
feat(api)!: exige autenticação em todos os endpoints

Antes, os endpoints de consulta eram públicos. Agora todo pedido
precisa do cabeçalho Authorization com um token válido.

BREAKING CHANGE: clientes sem token passam a receber 401.
Refs: #42
\end{terminalt}

\subsection{Os tipos}

A especificação só define dois tipos: \textbf{\arq{feat}} (novo recurso) e \textbf{\arq{fix}} (correção de bug). Os outros são convenção, herdada das regras do projeto Angular. A lista abaixo é a aceita pela configuração padrão do \emph{commitlint} (seção \ref{sec:commitlint}):

\begin{center}
\renewcommand{\arraystretch}{1.3}
\begin{tabularx}{\linewidth}{@{}>{\ttfamily\color{azulMB}}l X@{}}
\toprule
\textrm{\textbf{Tipo}} & \textbf{Quando usar} \\
\midrule
feat & um recurso novo para quem usa o sistema \\
fix & a correção de um bug \\
docs & só documentação (README, comentários, manuais) \\
style & formatação que não muda o comportamento (espaços, ponto e vírgula, indentação) \\
refactor & reorganização do código sem mudar o que ele faz nem corrigir bug \\
perf & melhoria de desempenho \\
test & criação ou ajuste de testes \\
build & sistema de build ou dependências (npm, pip, Maven\ldots) \\
ci & configuração de integração contínua (pipelines do GitLab, GitHub Actions) \\
chore & tarefas de manutenção que não se encaixam nos outros tipos \\
revert & desfaz um commit anterior \\
\bottomrule
\end{tabularx}
\end{center}

\begin{dica}[Tipos e versões]
Com Conventional Commits, ferramentas conseguem calcular a próxima versão pelo \emph{Versionamento Semântico}: \arq{fix} sobe o número de correção (1.2.\textbf{3}), \arq{feat} sobe o número de recurso (1.\textbf{3}.0) e um \arq{BREAKING CHANGE} sobe o número principal (\textbf{2}.0.0).
\end{dica}

\begin{atencao}[Tipos fora da lista]
Algumas referências populares em português acrescentam tipos como \arq{raw}, \arq{cleanup}, \arq{remove} e \arq{init}. Eles não fazem parte da especificação nem da configuração padrão do commitlint, que os rejeita. Só use se a equipe combinar e configurar isso.
\end{atencao}

\section{Boas práticas para a mensagem}

\begin{itemize}
  \item \textbf{Uma mudança por commit.} Se a descrição precisa de um \enquote{e}, provavelmente são dois commits. Commits pequenos são mais fáceis de revisar e de desfazer.
  \item \textbf{Descrição curta.} O livro \emph{Pro Git} recomenda uma primeira linha de até uns 50 caracteres e o corpo quebrado em 72 colunas. A configuração padrão do commitlint aceita até 100 caracteres no cabeçalho.
  \item \textbf{O corpo explica o porquê.} O \enquote{o quê} está no \cmd{git diff}. O motivo da mudança só existe na mensagem.
  \item \textbf{Um único modo verbal.} As convenções em inglês usam o imperativo (\enquote{add}, \enquote{fix}). Em português, equipes usam o imperativo (\enquote{corrige}), o presente (\enquote{corrige}) ou o gerúndio (\enquote{corrigindo}). A especificação não exige nenhum. Escolha um e mantenha.
  \item \textbf{Idioma combinado.} A especificação também não fixa idioma. Misturar português e inglês no mesmo histórico atrapalha a busca.
\end{itemize}

\section{Gitmoji: emojis nas mensagens}

O \termo{gitmoji} (\url{https://gitmoji.dev}) associa um emoji a cada tipo de mudança. No texto do commit, o emoji aparece como um código entre dois-pontos, o \emph{shortcode}. Alguns dos mais usados, com a descrição oficial traduzida:

\begin{center}
\renewcommand{\arraystretch}{1.25}
\small
\begin{tabularx}{\linewidth}{@{}>{\ttfamily\color{azulMB}}l X >{\ttfamily\color{azulMB}}l X@{}}
\toprule
\textrm{\textbf{Código}} & \textbf{Significado} & \textrm{\textbf{Código}} & \textbf{Significado} \\
\midrule
:sparkles: & novo recurso & :bug: & corrige um bug \\
:memo: & documentação & :art: & estrutura ou formato do código \\
:zap: & desempenho & :recycle: & refatoração \\
:fire: & remove código ou arquivos & :truck: & move ou renomeia arquivos \\
:white\_check\_mark: & testes & :construction\_worker: & CI \\
:wrench: & arquivos de configuração & :package: & arquivos compilados ou pacotes \\
:lock: & segurança ou privacidade & :ambulance: & correção crítica urgente \\
:pencil2: & corrige erro de digitação & :rewind: & reverte mudanças \\
:tada: & começa um projeto & :bookmark: & versão ou tag \\
:page\_facing\_up: & licença & :card\_file\_box: & mudanças de banco de dados \\
\bottomrule
\end{tabularx}
\end{center}

A lista completa tem 75 emojis e está em \url{https://gitmoji.dev}.

\begin{atencao}[Três cuidados com emojis]
\begin{enumerate}
  \item \textbf{O terminal não converte.} O GitHub mostra o emoji na página do commit, mas o \cmd{git log} exibe o texto \arq{:sparkles:} como foi escrito.
  \item \textbf{O mesmo emoji muda de sentido entre fontes.} Na lista oficial, \arq{:boom:} é quebra de compatibilidade, mas há tutoriais que o usam para \enquote{reverter}. Siga a lista oficial.
  \item \textbf{Emoji antes do tipo quebra o padrão.} A especificação exige que a mensagem comece pelo tipo. Uma mensagem como \arq{:bug: fix: corrige o login} é rejeitada pelo commitlint na configuração padrão.
\end{enumerate}
\end{atencao}

\section{Um padrão combinando os dois}

Muitas equipes juntam gitmoji e Conventional Commits num padrão próprio. É o caso do repositório pessoal do autor desta apostila, que usa:

\begin{terminalt}{· PADRÃO DE UMA EQUIPE}
<:emoji:> <tipo>(<escopo>): <Descrição no gerúndio>

:memo: docs(capacitacao): Adicionando os exercícios de assinatura
:bug: fix(login): Corrigindo o acesso com senha vazia
:truck: chore(docs): Movendo o manual para a pasta materials
\end{terminalt}

Não há nada de errado nisso, desde que \textbf{seja uma decisão da equipe, esteja escrita} (no README ou num guia de contribuição) e \textbf{as ferramentas estejam configuradas para aceitá-la}. O ponto é a consistência.

\section{Ferramentas}
\label{sec:commitlint}

\subsection*{commitlint: barrar mensagens fora do padrão}

O \termo{commitlint} (\url{https://commitlint.js.org}) confere cada mensagem antes de o commit ser gravado. Num projeto com Node.js:

\begin{terminal}
npm install -D @commitlint/cli @commitlint/config-conventional
echo "export default { extends: ['@commitlint/config-conventional'] };" > commitlint.config.js

# testar uma mensagem sem precisar fazer commit
echo "feat: adiciona exportação em PDF" | npx commitlint
\end{terminal}

Para rodar automaticamente em todo \cmd{git commit}, o guia oficial usa o \emph{husky}:

\begin{terminal}
npm install --save-dev husky
npx husky init
echo "npx --no -- commitlint --edit \$1" > .husky/commit-msg
\end{terminal}

Depois disso, \cmd{git commit -m "foo: teste"} é recusado, porque \arq{foo} não é um tipo válido.

\begin{atencao}[A configuração padrão é rígida]
Além de exigir o tipo no início, a configuração padrão rejeita descrições que começam com letra maiúscula. Um padrão como o da seção anterior precisa de regras próprias no \arq{commitlint.config.js}.
\end{atencao}

\subsection*{gitmoji-cli: escolher o emoji num menu}

\begin{terminal}
npm i -g gitmoji-cli
gitmoji -c        # monta o commit de forma interativa
gitmoji -l        # lista os emojis
\end{terminal}

\subsection*{Sem ferramenta nenhuma}

Nada disso é obrigatório. Um arquivo \arq{CONTRIBUTING.md} com o padrão e três exemplos já resolve boa parte do problema, e a revisão do Merge Request (capítulo \ref{cap:gitlab}) é um bom momento para lembrar o combinado.

\chapter{Assinatura de commits com GPG e SSH}
\label{cap:assinatura}

\section{O problema: o autor de um commit é só um texto}

No capítulo \ref{cap:antes} configuramos o nome e o e-mail com \cmd{git config}. O Git aceita \textbf{qualquer} valor ali, sem conferir nada. Veja:

\begin{terminal}
git -c user.name="Linus Torvalds" -c user.email="torvalds@linux-foundation.org" \
    commit --allow-empty -m "Commit que parece do Linus"
git log --format='%an <%ae>' -1
\end{terminal}
\begin{saida}
Linus Torvalds <torvalds@linux-foundation.org>
\end{saida}

O commit foi feito por nós, mas o histórico diz que foi o criador do Git. Se esse commit for enviado ao GitHub, ele aparece com o nome e a foto de quem tem aquele e-mail cadastrado.

\begin{conceito}[Assinatura de commit]
Assinar um commit é anexar a ele uma \textbf{assinatura criptográfica} feita com uma chave privada que só você tem. Quem tem a sua chave \textbf{pública} consegue conferir que o commit veio de você e que não foi alterado depois. O GitHub e o GitLab fazem essa conferência e mostram um selo ao lado do commit.
\end{conceito}

\subsection*{Os selos}

\begin{center}
\renewcommand{\arraystretch}{1.35}
\begin{tabularx}{\linewidth}{@{}>{\bfseries}l X@{}}
\toprule
\textbf{Selo} & \textbf{Significado} \\
\midrule
\color{verdeMB}Verified & o commit está assinado e a assinatura foi conferida com uma chave cadastrada na conta \\
\color{vermelhoAlerta}Unverified & o commit está assinado, mas a assinatura não pôde ser conferida (chave não cadastrada, e-mail diferente, chave de outra pessoa) \\
\textrm{(sem selo)} & o commit não foi assinado \\
\bottomrule
\end{tabularx}
\end{center}

\begin{atencao}[O e-mail precisa bater]
Para o selo \textbf{Verified}, o e-mail do commit (\cmd{git config user.email}) precisa ser um e-mail \textbf{verificado} na sua conta do GitHub ou do GitLab. No caso do GPG, ele também precisa estar dentro da chave. Essa é a causa mais comum de \enquote{Unverified}.
\end{atencao}

\section{GPG ou SSH?}

Há dois jeitos de assinar, e os dois são aceitos pelo GitHub e pelo GitLab:

\begin{center}
\renewcommand{\arraystretch}{1.35}
\small
\begin{tabularx}{\linewidth}{@{}>{\bfseries\color{azulMB}}l >{\raggedright\arraybackslash}X >{\raggedright\arraybackslash}X@{}}
\toprule
 & \textbf{GPG} & \textbf{SSH} \\
\midrule
Ferramenta & GnuPG (\cmd{gpg}) & OpenSSH (\cmd{ssh-keygen}), que já vem com o Git \\
Requisito & GnuPG instalado & Git 2.34 ou mais novo \\
Chave & uma chave própria, que carrega nome e e-mail & a mesma chave SSH que você usa para acessar o repositório pode servir \\
Validade & pode ter data de expiração embutida & a validade é definida no cadastro do GitLab \\
Conferir no terminal & direto, com o chaveiro do GPG & precisa de um arquivo \arq{allowed\_signers} \\
Quando escolher & padrão consolidado, usado há muito tempo & mais simples para quem já usa SSH \\
\bottomrule
\end{tabularx}
\end{center}

\begin{dica}
Escolha \textbf{um} dos dois e siga só a seção correspondente. Se não souber qual, comece pelo SSH: são menos passos.
\end{dica}

% -----------------------------------------------------------------------------
\section{Assinando com GPG}

\subsection{Instalar o GnuPG}

\begin{itemize}
  \item \textbf{Linux:} costuma vir instalado. Confira com \cmd{gpg --version}. No Debian e no Ubuntu, se faltar: \cmd{sudo apt install gnupg}.
  \item \textbf{Windows:} o Git for Windows já traz um \cmd{gpg} dentro do \textbf{Git Bash}. Use o Git Bash para todos os comandos desta seção.
\end{itemize}

\begin{atencao}[Windows: um GPG só]
Se você também instalar o Gpg4win, passa a ter \textbf{dois} GPGs no computador, cada um com o próprio chaveiro. Uma chave criada num não aparece no outro, e o Git reclama de \enquote{No secret key}. Crie a chave e assine sempre pelo mesmo. Para saber qual o Git está usando, rode \cmd{which gpg} no Git Bash.
\end{atencao}

\subsection{Passo 1 — Ver se você já tem uma chave}

\begin{terminal}
gpg --list-secret-keys --keyid-format=long
\end{terminal}

Se aparecer uma chave com o seu e-mail, pule para o passo 3.

\subsection{Passo 2 — Gerar a chave}

\begin{terminal}
gpg --full-generate-key
\end{terminal}

O comando faz algumas perguntas:

\begin{center}
\renewcommand{\arraystretch}{1.35}
\small
\begin{tabularx}{\linewidth}{@{}>{\bfseries}l >{\raggedright\arraybackslash}X@{}}
\toprule
\textbf{Pergunta} & \textbf{O que responder} \\
\midrule
Tipo de chave & aperte \tecla{Enter} para aceitar o padrão. Nas versões novas é ECC (ed25519); nas antigas, \enquote{RSA e RSA}. O GitHub e o GitLab aceitam os dois \\
Tamanho & só aparece para RSA: digite \texttt{4096}, como recomenda o GitLab \\
Validade & \texttt{0} = não expira, que é o padrão recomendado pelo GitHub. Uma data (como \texttt{1y}, um ano) é mais segura, mas exige renovar a chave \\
Nome e e-mail & o seu nome e \textbf{o mesmo e-mail verificado} no GitHub ou no GitLab \\
Senha & uma senha forte. Ela protege a chave se alguém copiar o seu computador \\
\bottomrule
\end{tabularx}
\end{center}

\subsection{Passo 3 — Descobrir o ID da chave}

\begin{terminal}
gpg --list-secret-keys --keyid-format=long
\end{terminal}
\begin{saida}
sec   rsa4096/30F2B65B9246B6CA 2017-08-18 [SC]
      D5E4F29F3275DC0CDA8FFC8730F2B65B9246B6CA
uid                   [ultimate] Maria Exemplo <maria@exemplo.com>
ssb   rsa4096/B7ABC0813E4028C0 2017-08-18 [E]
\end{saida}

O ID é o que vem \textbf{depois da barra} na linha \arq{sec}: aqui, \arq{30F2B65B9246B6CA}. Numa chave ECC, a linha começa com \arq{ed25519/}. O seu ID será outro.

\subsection{Passo 4 — Exportar a chave pública}

\begin{terminal}
gpg --armor --export 30F2B65B9246B6CA
\end{terminal}

Copie \textbf{tudo}, da linha \arq{-----BEGIN PGP PUBLIC KEY BLOCK-----} até a linha \arq{-----END PGP PUBLIC KEY BLOCK-----}, incluindo as duas.

\begin{atencao}[Pública, nunca privada]
O que vai para o site é a chave \textbf{pública} (\cmd{--export}). Nunca use \cmd{--export-secret-keys} para isso. A chave privada não sai do seu computador.
\end{atencao}

\subsection{Passo 5 — Cadastrar no GitHub}

\begin{enumerate}
  \item Clique na sua foto (canto superior direito) \faChevronRight\ \textbf{Settings}.
  \item Na seção \textbf{Access}, clique em \textbf{SSH and GPG keys}.
  \item Clique em \textbf{New GPG key}.
  \item Em \textbf{Title}, dê um nome (por exemplo, \enquote{Notebook pessoal}). Em \textbf{Key}, cole o bloco copiado.
  \item Clique em \textbf{Add GPG key}. O GitHub pode pedir a sua senha de novo.
\end{enumerate}

\subsection{Passo 5 (alternativo) — Cadastrar no GitLab}

\begin{enumerate}
  \item Clique no seu avatar (canto superior direito) \faChevronRight\ \textbf{Edit profile}.
  \item Na barra lateral, em \textbf{Access}, clique em \textbf{GPG keys}.
  \item Clique em \textbf{Add new key}, cole o bloco em \textbf{Key} e clique em \textbf{Add key}.
\end{enumerate}

Depois de cadastrada, a chave não pode ser editada: para mudar, remova e cadastre de novo. Num GitLab instalado pela instituição, o caminho é o mesmo.

\subsection{Passo 6 — Dizer ao Git qual chave usar}

\begin{terminal}
git config --global user.signingkey 30F2B65B9246B6CA
git config --global commit.gpgsign true     # assina todo commit
git config --global tag.gpgSign true        # assina toda tag
\end{terminal}

\begin{dica}[Já usou assinatura SSH antes?]
Se algum dia configurou \cmd{gpg.format ssh}, desfaça com \cmd{git config --global --unset gpg.format}. Sem isso, o Git continua tentando assinar com SSH.
\end{dica}

\subsection{Passo 7 — Ensinar o GPG a pedir a senha no terminal}

No Linux e no Git Bash, o GPG precisa saber em qual terminal pedir a senha da chave. Rode uma vez:

\begin{terminal}
[ -f ~/.bashrc ] && echo -e '\nexport GPG_TTY=$(tty)' >> ~/.bashrc
source ~/.bashrc
\end{terminal}

Quem usa \textbf{zsh} coloca a linha \arq{export GPG\_TTY=\$(tty)} no \arq{\textasciitilde/.zshrc}.

\subsection{Passo 8 — Testar}

\begin{terminal}
git commit --allow-empty -m "Testando a assinatura"
git log --show-signature -1
\end{terminal}
\begin{saida}
commit e55fbf5be5b86eb89737adea712d8c605b5f9529
gpg: Assinatura feita seg 28 set 2026 14:51:16 -03
gpg:                usando RSA chave 0D9FFCDB3C4EBDC3E7225C8CB4BEA86F1F6C7956
gpg: Assinatura correta de "Maria Exemplo <maria@exemplo.com>" [final]
Author: Maria Exemplo <maria@exemplo.com>
\end{saida}

A linha \enquote{Assinatura correta} (em inglês, \enquote{Good signature}) confirma que funcionou. Se você não ligou o \cmd{commit.gpgsign}, assine um commit específico com \cmd{git commit -S -m "..."}.

% -----------------------------------------------------------------------------
\section{Assinando com SSH}

\subsection{Requisitos}

\begin{itemize}
  \item \textbf{Git 2.34 ou mais novo} (\cmd{git --version});
  \item \textbf{OpenSSH 8.1 ou mais novo} (\cmd{ssh -V}). A versão 8.7 tem a assinatura quebrada: nela, atualize para a 8.8 ou mais nova;
  \item no GitLab, versão 15.7 ou mais nova. Num GitLab da instituição, confirme a versão com o administrador.
\end{itemize}

\subsection{Passo 1 — Ter uma chave SSH}

Se você já acessa o GitHub ou o GitLab por SSH, já tem uma: normalmente \arq{\textasciitilde/.ssh/id\_ed25519.pub}. Se não tiver, crie:

\begin{terminal}
ssh-keygen -t ed25519 -C "maria@exemplo.com"
\end{terminal}

Aperte \tecla{Enter} para aceitar o local padrão e escolha uma senha. O arquivo terminado em \arq{.pub} é a chave \textbf{pública}. O outro, sem extensão, é a \textbf{privada}, e nunca sai do seu computador.

\subsection{Passo 2 — Dizer ao Git para assinar com SSH}

\begin{terminal}
git config --global gpg.format ssh
git config --global user.signingkey ~/.ssh/id_ed25519.pub
git config --global commit.gpgsign true
git config --global tag.gpgSign true
\end{terminal}

\subsection{Passo 3 — Cadastrar como chave de assinatura no GitHub}

\begin{enumerate}
  \item Foto \faChevronRight\ \textbf{Settings} \faChevronRight\ \textbf{SSH and GPG keys} \faChevronRight\ \textbf{New SSH key}.
  \item Em \textbf{Title}, dê um nome. Em \textbf{Key type}, escolha \textbf{Signing Key}.
  \item Em \textbf{Key}, cole o conteúdo do arquivo \arq{.pub} (\cmd{cat \textasciitilde/.ssh/id\_ed25519.pub}) e clique em \textbf{Add SSH key}.
\end{enumerate}

\begin{atencao}[No GitHub, a mesma chave entra duas vezes]
Uma chave cadastrada como \emph{Authentication Key} serve para \cmd{git push}, mas \textbf{não} para verificar assinaturas. Se quiser usar a mesma chave para as duas coisas, cadastre-a duas vezes: uma com cada tipo.
\end{atencao}

Pelo terminal, com o GitHub CLI, dá no mesmo. A opção \cmd{--type} existe nas versões do \cmd{gh} lançadas a partir de 2023. Se a sua recusar a opção, use a página web.

\begin{terminal}
gh ssh-key add ~/.ssh/id_ed25519.pub --type signing --title "Notebook pessoal"
\end{terminal}

\subsection{Passo 3 (alternativo) — Cadastrar no GitLab}

\begin{enumerate}
  \item Avatar \faChevronRight\ \textbf{Edit profile} \faChevronRight\ \textbf{Access} \faChevronRight\ \textbf{SSH keys} \faChevronRight\ \textbf{Add new key}.
  \item Cole a chave pública em \textbf{Key} e dê um \textbf{Title}.
  \item Em \textbf{Usage type}, deixe \textbf{Authentication \& Signing} (o padrão) ou escolha \textbf{Signing}. Com \textbf{Authentication} sozinho, os commits ficam \enquote{Unverified}.
  \item Opcionalmente, defina uma \textbf{Expiration date} e clique em \textbf{Add key}.
\end{enumerate}

\subsection{Passo 4 — Permitir a conferência no seu computador}

O GitHub e o GitLab conferem a assinatura sozinhos. Mas, para o \emph{seu} Git conferir, ele precisa de uma lista de chaves confiáveis. Sem ela, aparece:

\begin{saida}
error: gpg.ssh.allowedSignersFile needs to be configured and exist for ssh signature verification
\end{saida}

Crie a lista uma vez:

\begin{terminal}
touch ~/.ssh/allowed_signers
git config --global gpg.ssh.allowedSignersFile ~/.ssh/allowed_signers
echo "$(git config --get user.email) namespaces=\"git\" $(cat ~/.ssh/id_ed25519.pub)" >> ~/.ssh/allowed_signers
\end{terminal}

Cada linha do arquivo associa um e-mail a uma chave. O \arq{namespaces="git"} limita a chave a assinaturas do Git.

\subsection{Passo 5 — Testar}

\begin{terminal}
git commit --allow-empty -m "Testando a assinatura SSH"
git log --show-signature -1
\end{terminal}
\begin{saida}
commit 86613eb51dee51c80a61f2cf8f1532df427e60c5
Good "git" signature for maria@exemplo.com with ED25519 key SHA256:Z/dW7mOFmUlo/Po1CSObYXoO3bAPY2bTohQtnJY0PGc
Author: Maria Exemplo <maria@exemplo.com>
\end{saida}

% -----------------------------------------------------------------------------
\section{Conferindo no GitHub e no GitLab}

Envie um commit assinado com \cmd{git push} e abra a lista de commits:

\begin{itemize}
  \item \textbf{GitHub:} na aba \textbf{Commits} do repositório, o selo \textbf{Verified} aparece ao lado do commit. Clique nele para ver qual chave assinou.
  \item \textbf{GitLab:} em \textbf{Code \faChevronRight\ Commits}, o selo \textbf{Verified} aparece da mesma forma.
\end{itemize}

No terminal, uma letra resume o estado de cada commit:

\begin{terminal}
git log --format='%h %G? %s'
\end{terminal}

\begin{center}
\renewcommand{\arraystretch}{1.3}
\begin{tabularx}{0.9\linewidth}{@{}>{\ttfamily\bfseries\color{azulMB}}l X@{}}
\toprule
G & assinatura válida \\
N & sem assinatura \\
B & assinatura inválida \\
E & não foi possível conferir (por exemplo, falta a chave pública ou o \arq{allowed\_signers}) \\
\bottomrule
\end{tabularx}
\end{center}

\begin{dica}[Vigilant mode no GitHub]
Em \textbf{Settings \faChevronRight\ SSH and GPG keys}, a opção \textbf{Vigilant mode} (\enquote{Flag unsigned commits as unverified}) faz o GitHub marcar como \textbf{Unverified} todo commit seu que não esteja assinado, inclusive os que alguém fizer se passando por você. Ligue só depois que todos os seus computadores estiverem assinando.
\end{dica}

% -----------------------------------------------------------------------------
\section{Cuidados com a chave}

\begin{itemize}
  \item \textbf{Guarde a chave privada e a senha.} Se perder a chave, é só criar outra e cadastrar, mas perde-se a continuidade.
  \item \textbf{Chave vazou? Revogue e cadastre outra.} Aqui as plataformas diferem:
  \begin{itemize}
    \item no \textbf{GitHub}, a verificação é registrada no momento do push. Os commits antigos \textbf{continuam Verified} mesmo que a chave seja removida, expire ou seja revogada;
    \item no \textbf{GitLab}, \textbf{Revoke} torna \textbf{Unverified} os commits antigos assinados com a chave, e \textbf{Remove} os mantém verificados. Use \emph{Revoke} quando a chave vazou.
  \end{itemize}
  \item \textbf{Commits feitos pela interface web} já saem assinados pela própria plataforma.
  \item No GitHub, um Pull Request integrado com \textbf{Rebase and merge} reescreve os commits, e eles perdem a sua assinatura.
  \item \textbf{Não é preciso publicar a chave em servidor de chaves} (\emph{keyserver}). Alguns tutoriais pedem isso, mas o GitHub e o GitLab só usam a chave cadastrada na conta.
\end{itemize}

\section{Outras plataformas e ferramentas}

\begin{itemize}
  \item \textbf{Bitbucket Cloud:} aceita GPG e SSH. A chave GPG entra em \textbf{Settings \faChevronRight\ Personal Bitbucket settings \faChevronRight\ GPG keys}. Para SSH, serve a mesma chave pessoal de acesso. Só commits enviados pela linha de comando são verificados.
  \item \textbf{1Password:} guarda a chave SSH no cofre e configura o Git sozinho (\arq{gpg.format}, \arq{user.signingkey}, \arq{commit.gpgsign} e o programa de assinatura dele). A chave continua precisando ser cadastrada como \emph{Signing Key} no GitHub ou com \emph{Usage type} de assinatura no GitLab.
\end{itemize}

\section{Problemas comuns}

\begin{center}
\renewcommand{\arraystretch}{1.4}
\small
\begin{tabularx}{\linewidth}{@{}>{\raggedright\arraybackslash\ttfamily\footnotesize}p{5.6cm} >{\raggedright\arraybackslash}X@{}}
\toprule
\textrm{\textbf{Mensagem ou sintoma}} & \textbf{O que fazer} \\
\midrule
error: gpg failed to sign the data & quase sempre é o \arq{GPG\_TTY} (passo 7 do GPG). Confira também se o \arq{user.signingkey} bate com \cmd{gpg --list-secret-keys} \\
gpg: signing failed: Inappropriate ioctl for device & falta o \arq{export GPG\_TTY=\$(tty)} \\
No secret key / secret key not available & o ID configurado não existe neste chaveiro. No Windows, confira se o Git usa o mesmo GPG em que a chave foi criada \\
gpg.ssh.allowedSignersFile needs to be configured and exist for ssh signature verification & crie o \arq{allowed\_signers} (passo 4 do SSH). A assinatura em si está certa \\
either user.signingkey or gpg.ssh.defaultKeyCommand needs to be configured & falta o \cmd{git config --global user.signingkey} \\
\textrm{Selo \textbf{Unverified} no site} & o e-mail do commit não está verificado na conta, a chave não foi cadastrada, ou a chave SSH foi cadastrada só para autenticação \\
\bottomrule
\end{tabularx}
\end{center}

\begin{dica}[Fontes oficiais]
Os passos deste capítulo seguem as documentações do GitHub (\url{https://docs.github.com/pt/authentication/managing-commit-signature-verification}) e do GitLab (\url{https://docs.gitlab.com/user/project/repository/signed_commits/}). Os nomes dos botões podem mudar com novas versões das plataformas. Se algo estiver diferente na sua tela, confira lá.
\end{dica}

\chapter{Problema 1: integrar um sistema já versionado ao GitLab institucional}
\label{cap:problema1}

Agora saímos da teoria. Este é um problema real, enfrentado no estágio.

\section{O cenário}

A instituição tem um GitLab empresarial restrito. Ao criar um repositório novo, ele já vem com:

\begin{itemize}
  \item uma branch \arq{main} criada;
  \item um \arq{README.md} padrão;
  \item regras institucionais;
  \item a \arq{main} \textbf{protegida}: não aceita \cmd{git push} direto, só Merge Request.
\end{itemize}

Só que vários sistemas \textbf{já existiam} nos computadores, e alguns já tinham anos de histórico Git próprio. Era preciso levar esses sistemas para o GitLab \textbf{sem perder o histórico}.

\section{Por que não funciona do jeito óbvio}

São dois históricos que nasceram separados:

\begin{center}
\begin{tikzpicture}
  \node[font=\small\bfseries\color{verdeMB}, anchor=west] at (-0.4,1.1) {REPOSITÓRIO LOCAL};
  \foreach \l/\x in {A/0,B/1.6,C/3.2,D/4.8}{\node[commit=verdeMB] (\l) at (\x,0.4) {\l};}
  \draw[hist=verdeMB] (A)--(B)--(C)--(D);
  \node[font=\small\bfseries\color{azulMB}, anchor=west] at (-0.4,-0.8) {REPOSITÓRIO GITLAB};
  \foreach \l/\x in {X/0,Y/1.6}{\node[commit] (\l) at (\x,-1.5) {\l};}
  \draw[hist] (X)--(Y);
  \node[font=\Large\color{vermelhoAlerta}] at (3.4,-1.5) {\faUnlink};
  \node[font=\small, anchor=west] at (4,-1.5) {nenhum ancestral em comum};
\end{tikzpicture}
\end{center}

O Git sabe juntar linhas que \textbf{nasceram de um mesmo ponto}. Entre \arq{A--B--C--D} e \arq{X--Y} não existe esse ponto. Por isso, um merge comum é recusado:

\begin{saida}
fatal: refusing to merge unrelated histories
\end{saida}

E um \cmd{git push} direto para a \arq{main} também é recusado, porque ela é protegida.

\section{O objetivo}

Transformar os dois históricos num só, sem perder nenhum commit:

\begin{center}
\begin{tikzpicture}
  \foreach \l/\x in {A/0,B/1.6,C/3.2,D/4.8}{\node[commit=verdeMB] (\l) at (\x,1) {\l};}
  \draw[hist=verdeMB] (A)--(B)--(C)--(D);
  \foreach \l/\x in {X/0,Y/2.4}{\node[commit] (\l) at (\x,-1) {\l};}
  \draw[hist] (X)--(Y);
  \node[commitm] (M) at (6.8,0) {M};
  \draw[hist=verdeMB] (D) to[out=0,in=135] (M);
  \draw[hist] (Y) to[out=0,in=-135] (M);
  \node[rotulo=verdeMB, anchor=west] at ([xshift=4mm]M.east) {desenvolvimento};
\end{tikzpicture}
\end{center}

O commit \textbf{M} é um commit de merge que une os dois históricos. Depois dele, a branch \arq{desenvolvimento} contém a \arq{main} institucional, então um Merge Request para a \arq{main} passa a funcionar normalmente.

\section{A solução, passo a passo}

\subsection*{Passo 1 — Adicionar o repositório remoto}

\begin{terminal}
git remote add origin <URL_DO_REPOSITORIO>
git remote -v
\end{terminal}
\begin{saida}
origin    <URL_DO_REPOSITORIO> (fetch)
origin    <URL_DO_REPOSITORIO> (push)
\end{saida}

Isso dá ao Git local um endereço remoto, chamado \arq{origin}. A URL aparece no botão \textbf{Code} (ou \textbf{Clone}) da página do projeto no GitLab.

\subsection*{Passo 2 — Buscar as informações do remoto}

\begin{terminal}
git fetch origin
\end{terminal}
\begin{saida}
From <URL_DO_REPOSITORIO>
 * [new branch]      main       -> origin/main
\end{saida}

O \cmd{fetch} baixa o histórico do GitLab para uma referência de leitura chamada \arq{origin/main}, \textbf{sem mexer} na sua branch nem nos seus arquivos.

\subsection*{Passo 3 — Renomear a branch local}

\begin{terminal}
git branch -m main desenvolvimento
\end{terminal}

A sua \arq{main} local vira \arq{desenvolvimento}. Isso evita confundi-la com a \arq{main} institucional, que é protegida e não pode receber push. Agora existem:

\begin{center}
\begin{tikzpicture}[l/.style={font=\small, anchor=west}]
  \node[rotulo=verdeMB] at (0,0) {desenvolvimento}; \node[l] at (1.8,0) {o histórico original do sistema (A--D)};
  \node[rotulo] at (-0.2,-0.8) {origin/main}; \node[l] at (1.8,-0.8) {o histórico institucional (X--Y)};
\end{tikzpicture}
\end{center}

\begin{dica}
Se a sua branch atual não se chama \arq{main}, use só \cmd{git branch -m desenvolvimento}, que renomeia a branch em que você está.
\end{dica}

\subsection*{Passo 4 — Unir os históricos}

\begin{terminal}
git merge --allow-unrelated-histories origin/main
\end{terminal}

O parâmetro \cmd{--allow-unrelated-histories} deve ser usado com consciência.

\begin{center}
\begin{tabularx}{\linewidth}{@{}XX@{}}
\begin{tcolorbox}[base, colback=vermelhoClaro, colframe=vermelhoAlerta, boxrule=0.8pt]
{\bfseries\color{vermelhoAlerta}\faTimes\; Ele NÃO significa}\\ \enquote{ignore todos os conflitos}.
\end{tcolorbox}
&
\begin{tcolorbox}[base, colback=verdeClaro, colframe=verdeMB, boxrule=0.8pt]
{\bfseries\color{verdeMB}\faCheck\; Ele significa}\\ \enquote{eu sei que esses históricos não têm ancestral comum e autorizo o Git a tentar integrá-los}.
\end{tcolorbox}
\end{tabularx}
\end{center}

Como os dois lados têm um \arq{README.md} diferente, surge um conflito:

\begin{saida}
Mesclagem automática de README.md
CONFLITO (adicionar/adicionar): conflito de mesclagem em README.md
Automatic merge failed; fix conflicts and then commit the result.
\end{saida}

\subsection*{Passo 5 — Resolver o conflito}

\begin{terminal}
git status
\end{terminal}
\begin{saida}
Caminhos não mesclados:
  (usar "git add <arquivo>..." para marcar resolução)
    ambos adicionaram:   README.md
\end{saida}

Aqui entra a parte que mais causa erro: \textbf{\texttt{ours} e \texttt{theirs}}.

\begin{center}
\begin{tikzpicture}
  \node[commitm, minimum size=13mm, font=\bfseries\color{azulMB}] (m) at (0,0) {merge};
  \node[caixa=verdeMB, text width=4.3cm] (o) at (-3.1,-2.2) {\textbf{OURS} (\enquote{nosso})\\ \small o lado da branch em que \textbf{você está} no momento do merge\\ \scriptsize aqui: \texttt{desenvolvimento}};
  \node[caixa, text width=4.3cm] (t) at (3.1,-2.2) {\textbf{THEIRS} (\enquote{deles})\\ \small o lado da branch que \textbf{está sendo incorporada}\\ \scriptsize aqui: \texttt{origin/main}};
  \draw[setafina] (m) -- (o); \draw[setafina] (m) -- (t);
\end{tikzpicture}
\end{center}

No nosso caso, queríamos manter o README do sistema, que está na branch em que estamos. Então:

\begin{terminal}
git restore --ours README.md
# forma tradicional, com o mesmo efeito:
git checkout --ours README.md
\end{terminal}

Se quiséssemos manter o README institucional:

\begin{terminal}
git restore --theirs README.md
# ou: git checkout --theirs README.md
\end{terminal}

\begin{atencao}[\texttt{ours} não é \enquote{o do meu computador}]
\texttt{ours} significa \textbf{o lado correspondente à branch atual durante aquele merge}. \texttt{theirs} é o outro lado. Os dois arquivos estão no seu computador. Interpretar errado faz você preservar a versão errada sem perceber.
\end{atencao}

\begin{atencao}[No rebase, os papéis se invertem]
Durante um \cmd{git rebase}, \texttt{ours} passa a ser a branch \emph{sobre a qual} você está reaplicando os commits, e \texttt{theirs} são os seus próprios commits. Nesta solução usamos \cmd{merge}, então vale a regra de cima. Se um dia usar rebase, confira antes.
\end{atencao}

\begin{dica}[E se eu quiser os dois READMEs?]
Você não precisa escolher um lado. Abra o \arq{README.md}, junte à mão as partes que interessam dos dois, apague os marcadores \texttt{<<<<<<<}, \texttt{=======} e \texttt{>>>>>>>}, e siga para o passo 6.
\end{dica}

\subsection*{Passo 6 — Registrar a resolução}

\begin{terminal}
git add README.md
git commit
\end{terminal}

Sem \cmd{-m}, o Git abre um editor com a mensagem \enquote{Merge remote-tracking branch 'origin/main' into desenvolvimento} já preenchida. Salve e feche. Este é o commit \textbf{M}: o histórico agora está integrado.

\begin{terminal}
git log --oneline --graph --all
\end{terminal}
\begin{saida}
*   89647ff Merge remote-tracking branch 'origin/main' into desenvolvimento
|\
| * 4738004 Initial commit
* 3ee51ff Ajuste
* 40d59e6 Primeira versão
\end{saida}

\subsection*{Passo 7 — Publicar a nova branch}

\begin{terminal}
git push origin desenvolvimento
\end{terminal}
\begin{saida}
 * [new branch]      desenvolvimento -> desenvolvimento
\end{saida}

No GitLab passam a existir, lado a lado, a \arq{main} (intocada e protegida) e a \arq{desenvolvimento} (o sistema com o histórico integrado).

\subsection*{Passo 8 — Abrir o Merge Request}

Este é o ponto final de todo o processo. Abra um Merge Request de \arq{desenvolvimento} para \arq{main}, pela interface web (capítulo \ref{cap:gitlab}) ou pelo terminal:

\begin{terminal}
glab mr create -s desenvolvimento -b main --title "Integra o sistema ao repositório institucional"
\end{terminal}

Como os dois históricos agora são relacionados e o conflito já foi tratado localmente, o merge \textbf{não é mais recusado}. A aprovação fica com quem tem permissão no repositório, conforme as regras do projeto.

\section{Resumo em oito linhas}

\begin{terminal}
git remote add origin <URL_DO_REPOSITORIO>          # 1. conectar
git fetch origin                                     # 2. baixar
git branch -m main desenvolvimento                   # 3. renomear
git merge --allow-unrelated-histories origin/main    # 4. unir
git restore --ours README.md                         # 5. decidir o conflito
git add README.md && git commit                      # 6. registrar
git push origin desenvolvimento                      # 7. publicar
glab mr create -s desenvolvimento -b main            # 8. pedir a integração
\end{terminal}

\section{Erros que podem aparecer}

\begin{center}
\renewcommand{\arraystretch}{1.4}
\small
\begin{tabularx}{\linewidth}{@{}>{\raggedright\arraybackslash\ttfamily\footnotesize}p{5.6cm} >{\raggedright\arraybackslash}X@{}}
\toprule
\textrm{\textbf{Mensagem}} & \textbf{O que fazer} \\
\midrule
fatal: refusing to merge unrelated histories & faltou o \cmd{--allow-unrelated-histories} no passo 4 \\
error: remote origin already exists. & o projeto já tem um \arq{origin}. Confira com \cmd{git remote -v}; para trocar a URL, use \cmd{git remote set-url origin <URL>} \\
remote: GitLab: You are not allowed to push code to protected branches on this project. & você tentou enviar para a \arq{main} protegida. Envie a \arq{desenvolvimento} e abra um Merge Request \\
error: Committing is not possible because you have unmerged files. & ainda há conflito sem resolver. Rode \cmd{git status}, resolva e faça \cmd{git add} \\
\bottomrule
\end{tabularx}
\end{center}

\begin{dica}
Quer treinar isso sem servidor nenhum? O exercício 4.1 (o desafio) do capítulo \ref{cap:exercicios} simula o GitLab institucional com uma pasta do seu computador e reproduz o caso inteiro.
\end{dica}

\chapter{Problema 2: documentar sistemas legados com um agente de IA}
\label{cap:problema2}

\section{Versionado não é documentado}

Resolvido o Problema 1, os sistemas passaram a ter histórico no GitLab. Mas surgiu outro problema:

\begin{conceito}[A ideia central]
\textbf{Um sistema estar versionado não significa que ele esteja documentado ou compreendido.}
\end{conceito}

Os sistemas existentes, alguns legados, em sua maioria não tinham documentação nenhuma. O quadro típico era este: documentação incompleta ou desatualizada, conhecimento concentrado em poucas pessoas, regras de negócio que só existiam no código, dependências desconhecidas, configurações pouco claras e manutenção difícil. Isso forma uma cadeia:

\begin{center}
\begin{tikzpicture}[p/.style={caixa, minimum width=2.1cm, text width=1.9cm, minimum height=1.15cm, font=\scriptsize}]
  \foreach \i/\t/\c in {0/Código\\existente/azulMB, 1/Pouca\\documentação/azulMB, 2/Conhecimento\\concentrado/azulMB,
                        3/Manutenção\\difícil/ouroMB!85!black, 4/Evolução\\difícil/ouroMB!85!black, 5/Risco\\operacional/vermelhoAlerta}{
    \node[p=\c] (n\i) at (\i*2.45,0) {\t};}
  \foreach \i [evaluate=\i as \j using int(\i+1)] in {0,...,4}{\draw[setafina] (n\i) -- (n\j);}
\end{tikzpicture}
\end{center}

\section{A abordagem: IA como ferramenta de engenharia}

A ferramenta usada para acelerar o trabalho foi um \termo{agente de IA para desenvolvimento}, o \textbf{Claude Code}, da Anthropic. Diferente de um chat comum, um agente desses trabalha dentro da pasta do projeto: lê os arquivos, roda comandos (com permissão) e escreve documentos.

\begin{atencao}[A ideia não é \enquote{mandar a IA escrever a documentação}]
A proposta é usar a IA como \textbf{ferramenta auxiliar de engenharia de software}: descobrir a arquitetura, analisar o código, identificar dependências, fluxos, regras de negócio, duplicações, problemas e riscos de segurança, sugerir refatorações e, por fim, produzir documentação.
\end{atencao}

\begin{center}
\begin{tikzpicture}[c/.style={caixa, minimum width=3.9cm, minimum height=2.9cm, text width=3.6cm, align=left, font=\small, anchor=north}]
  \node[c] (e1) at (0,0) {{\bfseries\color{azulMB}\faUserTie\; Engenheiro}\\[3pt] • define o objetivo\\ • define restrições\\ • orienta o agente};
  \node[c=ouroMB!85!black] (ia) at (5,0) {{\bfseries\color{azulMB}\includegraphics[height=3.5mm]{\imagens/claude-simbolo.pdf}\; Agente de IA}\\[3pt] • analisa\\ • sugere\\ • documenta\\ • executa o autorizado};
  \node[c=verdeMB] (e2) at (10,0) {{\bfseries\color{azulMB}\faUserCheck\; Engenheiro}\\[3pt] • revisa\\ • valida no sistema real\\ • decide};
  \draw[seta] (e1.east|-0,-1.45) -- (ia.west|-0,-1.45);
  \draw[seta] (ia.east|-0,-1.45) -- (e2.west|-0,-1.45);
\end{tikzpicture}
\end{center}

\section{A IA não substitui a validação humana}

O agente não é autoridade sobre o sistema. Ele pode interpretar errado regras de negócio, a intenção de quem escreveu o código, comportamentos implícitos, configurações, dependências e funcionalidades. Por isso o fluxo é sempre:

\begin{center}
\begin{tikzpicture}[p/.style={caixa, minimum width=2.5cm, text width=2.3cm, minimum height=1.1cm, font=\footnotesize}]
  \foreach \i/\t/\c in {0/Sistema/azulMB, 1/IA analisa/ouroMB!85!black, 2/IA produz hipóteses e documentação/ouroMB!85!black, 3/Humano valida/verdeMB, 4/Documentação consolidada/verdeMB}{
    \node[p=\c] (n\i) at (\i*3,0) {\t};}
  \foreach \i [evaluate=\i as \j using int(\i+1)] in {0,...,3}{\draw[setafina] (n\i) -- (n\j);}
\end{tikzpicture}
\end{center}

\begin{conceito}[Regra]
\textbf{A documentação gerada por IA deve ser validada contra o sistema real.} A responsabilidade final sobre as decisões técnicas continua sendo humana.
\end{conceito}

\section{A evolução real dos prompts}

\subsection{De onde vêm estes dados}

Um \termo{prompt} é a instrução que se dá ao agente. A qualidade do resultado depende muito dela. Para esta apostila, foi analisado o histórico de prompts do Claude Code do autor (o arquivo \arq{\textasciitilde/.claude/history.jsonl}), que registra cada pedido feito ao agente:

\begin{center}
\renewcommand{\arraystretch}{1.3}
\begin{tabular}{@{}>{\bfseries\color{azulMB}}r l@{}}
717 & prompts registrados \\
123 & sessões de trabalho \\
10/07 a 28/09/2026 & período coberto \\
8 & sistemas com trabalho de documentação completa (13/07 a 15/09) \\
\end{tabular}
\end{center}

Os nomes dos sistemas institucionais foram omitidos. Onde for preciso distingui-los, eles aparecem como Sistema A, B, C\ldots

\begin{atencao}[Limite dos dados]
O histórico guarda sempre o texto digitado, mas o texto dos prompts \emph{colados} só ficou salvo a partir de 03/09. Dos prompts colados antes disso restam só o tamanho e a data. As respostas do agente não foram analisadas, só os pedidos.
\end{atencao}

\subsection{Cinco versões}

\begin{center}
\renewcommand{\arraystretch}{1.4}
\footnotesize
\begin{tabularx}{\linewidth}{@{}>{\bfseries\color{azulMB}}l l l r >{\raggedright\arraybackslash}X@{}}
\toprule
\textbf{Versão} & \textbf{Início} & \textbf{Idioma} & \textbf{Tamanho} & \textbf{Características} \\
\midrule
V0 & 13/07 & PT & 953 & Sequência de 9 prompts colados: engenharia reversa, arquitetura, especificações, tarefas. Já pedia para \textbf{não modificar} e descobrir antes de escrever \\
V1 & 16/07 & PT & 231–944 & Pedidos curtos digitados, que cresciam a cada reuso. Em 24/07 surge o \textbf{GAP em vez de inferir}; em 27/07, \textbf{\enquote{não altere nenhuma linha de código}} \\
V2 & 10/08 & PT & 4.551 & Lista numerada de 47 itens: arquitetura, defeitos, LGPD (RIPD/DPIA), ADRs, UML, Docker, guia para leigos, projeto de evolução nos moldes do PMI \\
V3 & 14/09 & EN & 61.556 & Protocolo de refatoração: \textbf{fase 0 somente leitura}, regra de não pressupor, política de commits. \textbf{Sem exigência de idioma} \\
V4 & 15/09 & EN & 37.883 & Prompt final (\arq{model/SYSTEM\_DOCUMENTATION.md}): 42 seções, \textbf{saída obrigatória em PT-BR}, nunca inventar, AS-IS/TO-BE/GAP, \emph{Safe Stop-and-Report}, fases com validação \\
\bottomrule
\end{tabularx}
\end{center}
{\footnotesize\color{azulMedio} Tamanho em caracteres. \enquote{Início} é a data do primeiro uso de cada versão.}

\begin{center}
\begin{tikzpicture}
\begin{axis}[
  xbar, width=12cm, height=5.8cm, bar width=10pt,
  xmode=log, log basis x=10, xmin=100, xmax=300000,
  xlabel={\footnotesize tamanho do prompt em caracteres (escala logarítmica)},
  symbolic y coords={V4,V3,V2,V1,V0}, ytick=data,
  yticklabels={{\bfseries V4 · 15/09},{V3 · 14/09},{V2 · 10/08},{V1 · 16/07},{V0 · 13/07}},
  yticklabel style={font=\footnotesize\color{azulMB}},
  xticklabel style={font=\scriptsize\color{azulMedio}},
  axis x line*=bottom, axis y line*=left, x axis line style={azulMedio!40}, y axis line style={draw=none},
  xmajorgrids, grid style={azulMedio!15}, major tick style={draw=none},
  nodes near coords, nodes near coords style={font=\footnotesize\color{cinzaTexto}, anchor=west},
  point meta=explicit symbolic, enlarge y limits=0.14]
  \addplot[fill=azulMB, draw=none] coordinates {
    (37883,V4) [37.883] (61556,V3) [61.556] (4551,V2) [4.551] (944,V1) [231–944] (953,V0) [953]};
\end{axis}
\end{tikzpicture}
\end{center}

Do primeiro pedido livre (231 caracteres) ao prompt final (37.883), o tamanho cresceu cerca de \textbf{164 vezes}. O V4 é \textbf{38\,\% menor} que o V3: ficou mais enxuto e, ao mesmo tempo, mais rigoroso.

\subsection{Trechos que marcam a evolução}

\begin{analogia}[V0 — 13/07, português]
\enquote{Você é um Arquiteto de Software Sênior. Sua missão NÃO é modificar este projeto. [\ldots] Faça uma engenharia reversa completa [\ldots] Não escreva ainda a documentação final. Primeiro gere um relatório de descoberta [\ldots] Se faltar contexto, continue investigando automaticamente.}\\[3pt]
\footnotesize Já tinha a ideia de descobrir antes de escrever. Mas mandava o agente \enquote{continuar sozinho} quando faltasse contexto, e ainda não proibia inventar.
\end{analogia}

\begin{analogia}[V1 — 24/07, português]
\enquote{Se alguma informação estiver faltando, ela deve ser registrada como uma lacuna (GAP) em vez de ser inferida.}\\[3pt]
\footnotesize Primeira vez que a falta de informação vira um registro explícito, e não um palpite.
\end{analogia}

\begin{analogia}[V3 — 14/09, inglês]
\enquote{You are NOT authorized to simply rewrite the entire system at once. [\ldots] OBSERVE → INVENTORY → UNDERSTAND}\\[3pt]
\footnotesize \enquote{Você NÃO está autorizado a simplesmente reescrever o sistema inteiro de uma vez. Observar → inventariar → entender.}
\end{analogia}

\begin{analogia}[V4 — 15/09, inglês]
\enquote{ALL CONTENT YOU PRODUCE MUST BE WRITTEN IN BRAZILIAN PORTUGUESE.}\\[2pt]
\enquote{Never transform an assumption into a fact.}\\[3pt]
\footnotesize \enquote{Todo o conteúdo que você produzir deve ser escrito em português brasileiro.} \enquote{Nunca transforme uma suposição em fato.}
\end{analogia}

\section{O que deu errado e o que cada falha ensinou}

O histórico mostra os tropeços com clareza. Cada regra importante do prompt final apareceu \textbf{logo depois} de uma falha observada:

\begin{center}
\renewcommand{\arraystretch}{1.45}
\footnotesize
\begin{tabularx}{\linewidth}{@{}>{\raggedright\arraybackslash}X r >{\raggedright\arraybackslash}X@{}}
\toprule
\textbf{Falha observada} & \textbf{Ocorrências} & \textbf{Lição que virou regra} \\
\midrule
O agente parou no meio ou foi preciso pedir \enquote{continue de onde parou} & \textasciitilde27 & Trabalho em fases menores, com artefato salvo ao fim de cada fase \\
Limite de uso ou de contexto atingido & 14 + 10 & Idem; prompt mais enxuto do V3 para o V4 \\
Prompt reenviado ou reexecutado & 4 & Instruções mais precisas desde o início \\
Documentação superficial ou incompleta (\enquote{quero uma documentação mais profunda}) & 4 & Análise profunda em fases e auditoria de completude \\
Resposta em inglês depois do prompt em inglês & 1 & \S1 do V4: saída obrigatória em PT-BR \\
Agente criou ou alterou algo que não devia & \textasciitilde5 & Fase 0 somente leitura e \emph{Safe Stop-and-Report} \\
Commits com coautoria de IA & 4 lembretes & Cláusula explícita no protocolo de refatoração (V3) \\
Complexidade desnecessária (\enquote{preciso de coisas óbvias}) & 2 & Pedir clareza e público leigo \\
\bottomrule
\end{tabularx}
\end{center}
{\footnotesize\color{azulMedio} \enquote{14 + 10} = 14 avisos de limite de uso e 10 compactações de contexto. Contagens aproximadas, feitas sobre o texto dos prompts.}

\begin{dica}
Não houve nenhuma queixa explícita de informação inventada. A regra \enquote{nunca inventar} entrou de forma \textbf{preventiva}: primeiro como \enquote{GAP em vez de inferir} (24/07), depois como princípio absoluto no V4.
\end{dica}

\section{A melhora, em números}

\subsection*{Como foi medido}

Cada tarefa de documentação foi contada uma vez e classificada pela versão de prompt usada. Considerou-se \textbf{retrabalho} qualquer um destes sinais: pedir para refazer ou aprofundar a documentação do mesmo sistema, reenviar o mesmo prompt, pedir \enquote{continue}, ou corrigir idioma, pasta ou documento faltando.

\begin{center}
\renewcommand{\arraystretch}{1.35}
\small
\begin{tabular}{@{}>{\bfseries\color{azulMB}}l l r r r@{}}
\toprule
\textbf{Fase} & \textbf{Tipo de prompt} & \textbf{Tarefas} & \textbf{Com retrabalho} & \textbf{Taxa} \\
\midrule
V0 & pipeline estruturado, PT & 1 & 0 & \textasciitilde0\,\% \\
V1 & pedidos livres, PT & 12 & 6 & \textasciitilde50\,\% \\
V2 & lista numerada, PT & 3 & 2 & \textasciitilde67\,\% \\
V3 & protocolo em inglês, sem regra de idioma & 2 & 1 & \textasciitilde50\,\% \\
V4 & prompt final, inglês + saída PT-BR & 1 & 0 & \textasciitilde0\,\% \\
\midrule
\multicolumn{2}{@{}l}{\textbf{Prompts livres e lista (V1 + V2)}} & 15 & 8 & \textbf{\textasciitilde53\,\%} \\
\multicolumn{2}{@{}l}{\textbf{Prompts com contrato de execução (V0 + V3 + V4)}} & 4 & 1 & \textbf{\textasciitilde25\,\%} \\
\bottomrule
\end{tabular}
\end{center}

\begin{center}
\begin{tikzpicture}
  \node[align=center] (a) at (0,0) {{\titulofonte\fontsize{48}{48}\selectfont\color{vermelhoAlerta}\textasciitilde53\%}\\[-2pt]\footnotesize retrabalho com prompts livres};
  \node[font=\Huge\color{ouroMB}] at (3.6,0.3) {\faLongArrowAltRight};
  \node[align=center] (b) at (7.2,0) {{\titulofonte\fontsize{48}{48}\selectfont\color{verdeMB}\textasciitilde25\%}\\[-2pt]\footnotesize retrabalho com contrato de execução};
\end{tikzpicture}
\end{center}

\begin{atencao}[Leia com cautela]
É uma \textbf{estimativa indicativa}, não uma prova estatística:
\begin{itemize}
  \item a amostra é pequena (19 tarefas no total);
  \item depois de 15/09 só houve uma documentação completa com o V4 registrada no histórico;
  \item parte das paradas veio de limite de uso da ferramenta, e não da qualidade do prompt;
  \item o V2 foi pior que o V1 provavelmente porque pedia muito mais (47 itens, PDFs, projeto PMI), o que levava ao limite de uso e a paradas no meio.
\end{itemize}
Os sinais qualitativos, porém, são claros: cada falha gerou uma regra, e as regras resolveram as falhas que as motivaram.
\end{atencao}

\section{O prompt final por dentro}

O prompt final (V4) está no arquivo \arq{model/SYSTEM\_DOCUMENTATION.md}, que acompanha este material. Ele tem 42 seções, em inglês. Os blocos mais importantes são estes.

\subsection{Contrato de execução}

O prompt começa definindo \textbf{papéis} para o agente: engenheiro full-stack sênior, DevOps, arquiteto de software, engenheiro de segurança, QA, analista de sistemas, redator técnico, engenheiro de banco de dados, SRE, gerente de projetos e analista de negócios. E também a postura de um \textbf{professor} que valoriza clareza, organização e construção gradual do entendimento.

A missão declarada \textbf{não} é \enquote{escrever documentação}. É descobrir, entender, fazer engenharia reversa, validar contra o código, identificar o comportamento real, documentar o estado atual, apontar defeitos e riscos, explicar para públicos técnicos e leigos, propor uma evolução justificada e, por fim, \textbf{validar de forma independente} o que foi produzido.

\subsection{Idioma obrigatório}

Tudo o que o agente produz (arquivos \arq{.md} e \arq{.pdf}, README, rótulos de diagramas, tabelas, manuais, planos e relatórios) deve estar em \textbf{português brasileiro} formal. As exceções são elementos técnicos que não se traduzem, como nomes de funções, variáveis e arquivos, comandos, bibliotecas e código-fonte. Mesmo nesses casos, a explicação em volta fica em português.

\subsection{Nunca inventar}

Todo fato sobre o sistema precisa de evidência: código, configuração, banco, migrações, testes, infraestrutura, documentação existente, logs. O que não puder ser confirmado recebe um rótulo explícito:

\begin{center}
\ttfamily\small\color{azulMB}
DESCONHECIDO \enspace·\enspace NÃO VERIFICADO \enspace·\enspace NÃO ENCONTRADO \enspace·\enspace INFERIDO\\
RECOMENDAÇÃO \enspace·\enspace PROPOSTA \enspace·\enspace REQUER VALIDAÇÃO
\end{center}

\subsection{Separação de estados}

A documentação nunca mistura:

\begin{center}
\renewcommand{\arraystretch}{1.3}
\begin{tabularx}{0.9\linewidth}{@{}>{\bfseries\color{azulMB}}l X@{}}
\toprule
AS-IS & como o sistema \textbf{é hoje} \\
TO-BE & como deveria ser depois de uma evolução proposta \\
GAP & a diferença entre os dois \\
RECOMMENDATION & uma melhoria tecnicamente justificada \\
PROPOSED IMPLEMENTATION & uma forma possível de implementar a recomendação \\
\bottomrule
\end{tabularx}
\end{center}

\subsection{Somente leitura na descoberta}

Durante a descoberta e a análise, o agente \textbf{não} refatora, não altera código, não apaga nem renomeia arquivos, não muda dependências, infraestrutura, banco ou configuração, e não executa operações destrutivas. Só pode criar documentação em \arq{docs/}.

\begin{conceito}[Descobrir antes de modificar]
Isso reduz o risco de a IA fazer alterações prematuras baseadas numa compreensão incompleta do sistema.
\end{conceito}

\subsection{Safe Stop-and-Report}

Um agente que trabalha sobre um sistema real precisa saber \textbf{quando parar}. O protocolo lista as situações em que o agente deve interromper a operação e reportar, por exemplo:

\begin{itemize}
  \item falta informação que o repositório não fornece;
  \item duas fontes de evidência se contradizem;
  \item seria preciso uma operação destrutiva ou afetar produção, banco real ou dados de usuários;
  \item o agente encontrou senhas, tokens, chaves ou dados pessoais;
  \item continuar exigiria inventar informação ou supor uma autorização que não foi dada;
  \item é preciso aprovação humana para seguir.
\end{itemize}

\begin{center}
\begin{tikzpicture}[p/.style={caixa, minimum width=3cm, minimum height=1cm, font=\small\bfseries}]
  \node[p=vermelhoAlerta] (a) at (0,0) {Problema crítico};
  \node[p] (b) at (3.6,0) {PARAR};
  \node[p] (c) at (7.2,0) {REPORTAR};
  \node[p=verdeMB] (d) at (10.8,0) {AGUARDAR};
  \draw[setafina] (a)--(b); \draw[setafina] (b)--(c); \draw[setafina] (c)--(d);
  \node[font=\scriptsize, below=1mm of d] {a decisão humana};
\end{tikzpicture}
\end{center}

Cada parada gera um relatório padronizado (\enquote{SAFE STOP REPORT}) com o que estava sendo feito, o motivo, a evidência, os fatos verificados, o risco, o que \emph{não} foi alterado e qual decisão é necessária. A gravidade vai de \textbf{STOP-P0} (crítico, parada imediata) a \textbf{STOP-P3} (baixo). E uma parada \textbf{não é fracasso}: é o comportamento correto diante do risco.

\subsection{Documentação em fases}

O prompt proíbe fazer tudo numa tacada só. O trabalho segue fases, cada uma com um artefato:

\begin{center}
\renewcommand{\arraystretch}{1.3}
\small
\begin{tabularx}{\linewidth}{@{}>{\bfseries\color{azulMB}}l >{\raggedright\arraybackslash}X >{\ttfamily\footnotesize}l@{}}
\toprule
\textbf{Fase} & \textbf{Objetivo} & \textrm{\textbf{Artefato}} \\
\midrule
0 & Descoberta somente leitura: estrutura, tecnologias, dependências, APIs, banco & docs/00-discovery/DISCOVERY.md \\
1 & Mapa do sistema: componentes, relações, entrada, processamento e saída de dados & docs/00-discovery/SYSTEM\_MAP.md \\
2 & Análise profunda: fluxos, regras, persistência, autenticação, integrações & \textrm{(alimenta as seguintes)} \\
3 & Modelo de conhecimento: entidades, regras, fluxos, riscos, lacunas & docs/00-discovery/KNOWLEDGE\_MODEL.md \\
4 & Validação do entendimento: contradições, conclusões sem evidência, esquecimentos & docs/00-discovery/VALIDATION.md \\
5+ & Documentação técnica, regras de negócio, dados, APIs, segurança, operação, README, guia para leigos, projeto de evolução e auditorias & docs/ \\
\bottomrule
\end{tabularx}
\end{center}

\begin{conceito}[Por que em fases]
Para evitar a cascata \enquote{entendimento errado \faLongArrowAltRight\ documentação errada \faLongArrowAltRight\ recomendação errada \faLongArrowAltRight\ projeto futuro errado}. Em vez disso: \textbf{descobrir \faLongArrowAltRight\ evidenciar \faLongArrowAltRight\ modelar \faLongArrowAltRight\ validar \faLongArrowAltRight\ documentar \faLongArrowAltRight\ auditar \faLongArrowAltRight\ evoluir}.
\end{conceito}

\subsection{O restante do prompt}

As demais seções pedem: visão geral, arquitetura e tecnologias; regras de negócio (atuais e recomendadas); arquitetura de dados; APIs; telas; segurança; LGPD (RIPD/DPIA); defeitos, falhas, fraquezas e dívida técnica; modelo de qualidade e testes; operação, Docker e recuperação de desastres; plano de manutenção e de desenvolvimento; modelo de ameaças; ADRs (registros de decisão de arquitetura); governança; glossário; diagramas; requisitos; README profissional; documentação para leigos (manual do usuário, guia completo, runbook); projeto de evolução alinhado ao PMI; priorização de P0 a P3; rastreabilidade; auditorias de qualidade e de completude; geração de PDF; e um relatório executivo final. O fecho resume o espírito:

\begin{analogia}[Princípio final do prompt]
\enquote{Your objective is not to make the system look better than it actually is. Your objective is to make the system \textbf{understood exactly as it is}.}\\[3pt]
\footnotesize \enquote{O objetivo não é fazer o sistema parecer melhor do que é. O objetivo é fazer o sistema ser \textbf{entendido exatamente como ele é}.}
\end{analogia}

\section{Prompt em inglês, documentação em português}

O prompt final foi mantido em inglês. A motivação adotada foi a \textbf{economia de tokens}: os modelos de IA cobram e limitam o trabalho por \emph{tokens} (pedaços de texto), e textos em inglês costumam ser divididos em menos tokens que o mesmo conteúdo em português.

\begin{atencao}[O que os dados mostram]
O histórico confirma a troca de idioma e o efeito colateral dela, mas \textbf{não mediu} a economia de tokens. Trate essa economia como uma prática adotada, não como resultado comprovado aqui.
\end{atencao}

O efeito colateral foi real: no dia 14/09, o prompt V3, só em inglês, fez o agente responder em inglês. No dia seguinte, o V4 ganhou a regra de idioma. A lição:

\begin{conceito}[Idioma do prompt × idioma do resultado]
São coisas diferentes. Se você escreve o prompt em inglês e quer o resultado em português, \textbf{diga isso explicitamente}.
\end{conceito}

\section{O que o prompt pede e o que ele entrega}

O prompt pode ser resumido em seis grandes perguntas:

\begin{center}
\renewcommand{\arraystretch}{1.35}
\begin{tabularx}{0.9\linewidth}{@{}>{\bfseries\color{verdeMB}}l X@{}}
\toprule
Descoberta & O que existe neste sistema? \\
Compreensão & Como as partes do sistema se relacionam? \\
Funcionamento & Como o sistema funciona? \\
Auditoria & Quais são os problemas, riscos e pontos fracos? \\
Documentação & Como transformar esse conhecimento em documentação útil? \\
Validação & Como verificar se a documentação corresponde ao sistema real? \\
\bottomrule
\end{tabularx}
\end{center}

O retorno esperado não é \enquote{vários arquivos Markdown}. É transformar \enquote{código que poucas pessoas conhecem} em \enquote{sistema cujo funcionamento pode ser compreendido e mantido por outras pessoas}. Uma pessoa nova deve conseguir, aos poucos: entender o propósito do sistema, instalar o ambiente, compreender a arquitetura e o banco, entender os principais fluxos, localizar funcionalidades, identificar dependências, diagnosticar problemas, fazer manutenção e continuar a evolução.

Um exemplo de estrutura produzida (a exata se adapta a cada sistema):

\begin{center}
\begin{minipage}{0.62\linewidth}
\dirtree{%
.1 docs/.
.2 00-discovery/.
.3 DISCOVERY.md.
.3 SYSTEM\_MAP.md.
.3 KNOWLEDGE\_MODEL.md.
.3 VALIDATION.md.
.2 architecture.md.
.2 business-rules.md.
.2 database.md.
.2 api.md.
.2 security.md.
.2 known-issues.md.
.2 technical-debt.md.
.2 guia-para-leigos.md.
}
\end{minipage}
\end{center}

\section{Como usar na prática}

\begin{atencao}[Autorização primeiro]
Antes de rodar um agente de IA sobre um sistema da instituição, confirme que o uso está autorizado e segue as normas de segurança da informação da sua OM. O agente lê o código e o envia a um serviço externo para ser processado.
\end{atencao}

\begin{enumerate}
  \item Instale o Claude Code seguindo a documentação oficial: \url{https://docs.claude.com/en/docs/claude-code/overview}.
  \item Garanta que o projeto está versionado e \textbf{limpo} (\cmd{git status} sem pendências). Assim, qualquer alteração do agente aparece no \cmd{git diff} e pode ser desfeita.
  \item Crie uma branch para a documentação: \cmd{git switch -c docs/documentacao-inicial}.
  \item Abra o terminal na pasta do projeto e inicie o agente com \cmd{claude}.
  \item Entregue o prompt: cole o conteúdo ou peça para o agente ler o arquivo do prompt.
  \item Acompanhe. Quando o agente parar e reportar (\emph{safe stop}), leia o relatório e decida.
  \item \textbf{Valide} a documentação contra o sistema real (veja a lista abaixo).
  \item Registre a documentação com commit, envie a branch e abra um Merge Request, para que a equipe revise.
\end{enumerate}

\subsection*{Lista de validação}

\begin{itemize}
  \item A arquitetura documentada corresponde ao código?
  \item Os endpoints citados realmente existem?
  \item O banco de dados corresponde ao que está escrito?
  \item As regras de negócio estão corretas? (pergunte a quem usa o sistema)
  \item As instruções de instalação funcionam numa máquina limpa?
  \item Os problemas apontados são reais?
  \item O que ficou marcado como \texttt{DESCONHECIDO} ou \texttt{REQUER VALIDAÇÃO} foi tratado?
  \item Nenhuma senha, token ou dado pessoal foi copiado para a documentação?
\end{itemize}

\begin{dica}
A documentação é um \textbf{artefato vivo}. Quando o sistema mudar, ela também precisa ser revisada. Coloque a atualização da documentação como item do próprio Merge Request.
\end{dica}

\chapter{Exercícios práticos}
\label{cap:exercicios}

Os exercícios seguem a ordem da apostila. As partes 1, 4 e 5 e o exercício 6.1 funcionam só com o Git instalado, sem internet e sem servidor. A parte 2 precisa de acesso a um GitLab. A parte 3 precisa de um agente de IA e de um projeto que você possa usar. Os exercícios 6.2 e 6.3 precisam de uma conta no GitHub ou no GitLab.

\begin{dica}
Faça os exercícios numa pasta de treino, fora de qualquer projeto real. Se algo der errado, apague a pasta e comece de novo: é para isso que ela existe.
\end{dica}

% -----------------------------------------------------------------------------
\section{Parte 1 — Git no seu computador}

\begin{exercicio}{1.1 — Primeiro repositório e primeiro commit}
\begin{enumerate}
  \item Crie um repositório local.
  \item Crie um arquivo.
  \item Execute \cmd{git add}.
  \item Execute \cmd{git commit}.
\end{enumerate}
\end{exercicio}

\begin{terminal}
mkdir exercicio-git
cd exercicio-git
git init -b main
echo "Olá, Git!" > saudacao.txt
git status
git add saudacao.txt
git commit -m "Cria a saudação"
\end{terminal}

\begin{verifique}
\cmd{git log --oneline} mostra uma linha com \enquote{Cria a saudação}, e \cmd{git status} diz que não há nada a submeter.
\end{verifique}

\begin{exercicio}{1.2 — Branch, alteração e merge}
\begin{enumerate}\setcounter{enumi}{4}
  \item Crie uma branch.
  \item Altere o arquivo.
  \item Crie outro commit.
  \item Faça o merge na \arq{main}.
\end{enumerate}
\end{exercicio}

\begin{terminal}
git switch -c feature/despedida
echo "Até logo!" >> saudacao.txt       # >> acrescenta uma linha no fim
git commit -am "Adiciona a despedida"  # -a prepara os arquivos já acompanhados
git switch main
git merge feature/despedida
cat saudacao.txt
\end{terminal}

\begin{verifique}
O merge diz \enquote{Fast-forward} e o arquivo tem duas linhas: \enquote{Olá, Git!} e \enquote{Até logo!}.
\end{verifique}

\begin{exercicio}{1.3 — Simular e resolver um conflito}
\begin{enumerate}\setcounter{enumi}{8}
  \item Simule um conflito.
  \item Resolva o conflito.
\end{enumerate}
A ideia é mudar \textbf{a mesma linha} de formas diferentes em duas branches.
\end{exercicio}

\begin{terminal}
# na branch nova, a saudação fica formal
git switch -c feature/formal
printf "Prezados, bom dia!\nAté logo!\n" > saudacao.txt
git commit -am "Deixa a saudação formal"

# na main, a mesma linha fica informal
git switch main
printf "E aí, pessoal!\nAté logo!\n" > saudacao.txt
git commit -am "Deixa a saudação informal"

# agora junte as duas
git merge feature/formal
\end{terminal}
\begin{saida}
Mesclagem automática de saudacao.txt
CONFLITO (conteúdo): conflito de mesclagem em saudacao.txt
Automatic merge failed; fix conflicts and then commit the result.
\end{saida}

Abra \arq{saudacao.txt} num editor de texto. Você vai ver os marcadores explicados na seção \ref{sec:conflitos}. Escolha a versão que fica (ou escreva uma terceira), apague os marcadores, salve e registre:

\begin{terminal}
git add saudacao.txt
git commit -m "Resolve o conflito da saudação"
git log --oneline --graph --all
\end{terminal}
\begin{saida}
*   757aae3 Resolve o conflito da saudação
|\
| * c2524a8 Deixa a saudação formal
* | a41aded Deixa a saudação informal
|/
* afe16bb Adiciona a despedida
* 4203776 Cria a saudação
\end{saida}

\begin{verifique}
O desenho do \cmd{git log --graph} mostra as duas linhas se separando e se juntando de novo no commit de resolução. Os códigos dos commits no seu computador serão outros.
\end{verifique}

% -----------------------------------------------------------------------------
\section{Parte 2 — GitLab}

\begin{exercicio}{2.1 — Um Merge Request completo}
\begin{enumerate}
  \item Crie ou use um projeto no GitLab.
  \item Crie uma branch.
  \item Faça uma alteração.
  \item Faça \cmd{push}.
  \item Crie um Merge Request.
  \item Revise a alteração.
  \item Faça o merge.
\end{enumerate}
\end{exercicio}

\begin{terminal}
git clone <URL_DO_PROJETO_DE_TREINO>
cd <pasta-do-projeto>
git switch -c exercicio/meu-nome
echo "Treinamento de Git na Marinha" >> README.md
git add README.md
git commit -m "Registra o treinamento no README"
git push origin exercicio/meu-nome
\end{terminal}

Depois, no navegador:

\begin{enumerate}
  \item abra o projeto e clique em \textbf{Create merge request} (ou vá em \textbf{Merge requests \faChevronRight\ New merge request});
  \item confira: origem \arq{exercicio/meu-nome}, destino \arq{main};
  \item escreva um título e uma descrição e crie o MR;
  \item abra a aba \textbf{Changes} e deixe um comentário numa linha;
  \item peça a um colega para aprovar e faça o merge (se o projeto permitir).
\end{enumerate}

\begin{dica}[Pelo terminal]
Os passos no navegador podem ser trocados por \cmd{glab mr create -b main}, \cmd{glab mr view} e \cmd{glab mr merge} (capítulo \ref{cap:glab}).
\end{dica}

\begin{verifique}
O MR aparece como \enquote{Merged}, e a linha nova está no \arq{README.md} da \arq{main}.
\end{verifique}

% -----------------------------------------------------------------------------
\section{Parte 3 — IA}

\begin{atencao}[Escolha bem o projeto]
Use um projeto pequeno, \textbf{pessoal ou de código aberto}. Não use código institucional sem autorização.
\end{atencao}

\begin{exercicio}{3.1 — Descoberta com um agente de IA}
Escolha um pequeno projeto existente e peça ao agente:
\begin{enumerate}
  \item realizar descoberta somente leitura;
  \item identificar tecnologias;
  \item identificar a arquitetura;
  \item identificar dependências;
  \item identificar funcionalidades;
  \item identificar problemas;
  \item gerar documentação;
  \item apresentar dúvidas e pontos que exigem validação humana.
\end{enumerate}
\end{exercicio}

Um ponto de partida é o prompt abaixo. Ele é uma \textbf{versão reduzida}, escrita para este exercício, que aproveita as travas do prompt final: somente leitura, nunca inventar, parar e perguntar.

\begin{terminalt}{· PROMPT DE EXEMPLO}
Você vai analisar este projeto. Siga estas regras sem exceção:

1. MODO SOMENTE LEITURA: não altere, não apague e não renomeie nenhum
   arquivo, e não faça commits. Você só pode criar arquivos em docs/.
2. NUNCA INVENTE: todo fato precisa de evidência (arquivo e linha).
   O que não puder confirmar, marque como DESCONHECIDO, INFERIDO ou
   REQUER VALIDAÇÃO.
3. PARE E PERGUNTE: se encontrar contradição, senha, token, dado pessoal
   ou algo que exija decisão humana, pare e me relate antes de seguir.
4. IDIOMA: escreva tudo em português brasileiro.

Tarefa:
a) Crie docs/DISCOVERY.md com a estrutura de pastas, as tecnologias e
   as dependências encontradas.
b) Crie docs/ARQUITETURA.md explicando os componentes e como se
   comunicam, com um diagrama em texto.
c) Crie docs/FUNCIONALIDADES.md listando o que o sistema faz.
d) Crie docs/PROBLEMAS.md com defeitos, riscos e dívida técnica.
e) Termine com uma lista de perguntas que eu preciso validar.
\end{terminalt}

\begin{exercicio}{3.2 — Validar o que a IA produziu}
Pegue três afirmações da documentação gerada e confira cada uma no código. Anote: estava certa, errada ou incompleta? Responda também às perguntas que o agente deixou no item (e).
\end{exercicio}

\begin{verifique}
\cmd{git status} mostra \textbf{só} arquivos novos em \arq{docs/}. Se aparecer qualquer outro arquivo modificado, o agente desrespeitou a regra 1: use \cmd{git diff} para ver o que mudou e \cmd{git restore} para desfazer.
\end{verifique}

% -----------------------------------------------------------------------------
\section{Parte 4 — Desafio: reproduzir o Problema 1}

\begin{exercicio}{4.1 — O caso real, sem servidor}
Simule o GitLab institucional com um repositório \enquote{vazio de trabalho} (\emph{bare}) numa pasta do seu computador, com a \arq{main} e o README padrão. Depois, integre a ele um sistema que já tem histórico próprio, seguindo os oito passos do capítulo \ref{cap:problema1}.
\end{exercicio}

\textbf{Preparação: o \enquote{GitLab} de mentira.}

\begin{terminal}
mkdir desafio && cd desafio
git init --bare -b main gitlab.git       # faz o papel do servidor
git clone gitlab.git instituicao
cd instituicao
echo "# Projeto (README padrão da instituição)" > README.md
git add README.md
git commit -m "Initial commit"
git push origin main
cd ..
\end{terminal}

\textbf{Preparação: o sistema que já existia.}

\begin{terminal}
mkdir sistema && cd sistema
git init -b main
echo "# Sistema de Controle" > README.md
echo "print('ola')" > app.py
git add README.md app.py
git commit -m "Primeira versão"
echo "x = 1" >> app.py
git commit -am "Ajuste"
\end{terminal}

\textbf{Agora é com você.} Ainda dentro de \arq{sistema/}, aplique os oito passos. Como o \enquote{servidor} é uma pasta, a URL é o caminho dela: \arq{../gitlab.git}. O passo 8 (Merge Request) não existe numa pasta, então pare no passo 7.

\begin{enumerate}
  \item Antes do passo 4, tente \cmd{git merge origin/main} sem o parâmetro. Qual mensagem aparece? Por quê?
  \item No passo 5, mantenha o README do \textbf{sistema}. Qual comando você usou: \cmd{--ours} ou \cmd{--theirs}?
  \item Depois do passo 7, rode \cmd{git log --oneline --graph --all}. Onde está o commit que une os dois históricos?
\end{enumerate}

\begin{tcolorbox}[base, colback=cinzaSuave, colframe=cinzaSuave, title={\faKey\enspace Respostas}, coltitle=azulMB, attach title to upper={\\[2pt]}]
\begin{enumerate}
  \item \texttt{fatal: refusing to merge unrelated histories}. Os dois repositórios nasceram separados e não têm ancestral comum.
  \item \cmd{git restore --ours README.md}: durante o merge você está na branch \arq{desenvolvimento}, que tem o README do sistema.
  \item No topo: o commit \enquote{Merge remote-tracking branch 'origin/main' into desenvolvimento}, com duas linhas chegando nele. Uma traz \enquote{Primeira versão} e \enquote{Ajuste}, a outra traz \enquote{Initial commit}.
\end{enumerate}
\end{tcolorbox}

% -----------------------------------------------------------------------------
\section{Parte 5 — Padrões de commit}

\begin{exercicio}{5.1 — Reescrever mensagens ruins}
Reescreva cada mensagem abaixo no formato Conventional Commits (capítulo \ref{cap:padroes}). Escolha o tipo, invente um escopo que faça sentido e escreva uma descrição curta.
\begin{enumerate}
  \item \enquote{arrumei o erro que dava quando a senha tava vazia}
  \item \enquote{coloquei explicação de como instalar no readme}
  \item \enquote{mudei a identação do arquivo de rotas}
  \item \enquote{agora o relatório sai em pdf também}
  \item \enquote{atualizei o pipeline pra rodar os testes}
  \item \enquote{a API agora exige token em tudo, quem não mandar leva 401}
\end{enumerate}
\end{exercicio}

\begin{exercicio}{5.2 — Aplicar no seu repositório de treino}
No repositório \arq{exercicio-git} da Parte 1, faça três commits pequenos, cada um com um tipo diferente (\arq{feat}, \arq{docs} e \arq{fix}). Depois rode \cmd{git log --oneline} e confira se dá para entender o histórico só lendo as mensagens.
\end{exercicio}

\begin{tcolorbox}[base, colback=cinzaSuave, colframe=cinzaSuave, title={\faKey\enspace Respostas possíveis do 5.1}, coltitle=azulMB, attach title to upper={\\[2pt]}]
\ttfamily\small
\begin{enumerate}
  \item fix(login): aceita tentativa com senha vazia sem erro
  \item docs(readme): explica a instalação
  \item style(rotas): corrige a indentação
  \item feat(relatorio): exporta em PDF
  \item ci: roda os testes no pipeline
  \item feat(api)!: exige token em todos os endpoints
\end{enumerate}
\rmfamily\normalsize
No item 6, o \arq{!} marca a quebra de compatibilidade. Um rodapé \arq{BREAKING CHANGE: clientes sem token recebem 401} deixaria isso ainda mais claro. Os escopos e as palavras podem variar: o importante é o tipo certo e uma descrição que diga o que mudou.
\end{tcolorbox}

% -----------------------------------------------------------------------------
\section{Parte 6 — Assinatura de commits}

\begin{exercicio}{6.1 — Ver o problema com os próprios olhos}
Num repositório de treino, faça um commit com um nome e um e-mail que não são seus e confira o que o \cmd{git log} mostra:
\end{exercicio}

\begin{terminal}
git -c user.name="Pessoa Qualquer" -c user.email="qualquer@exemplo.com" \
    commit --allow-empty -m "Quem fez este commit?"
git log --format='%an <%ae> | assinatura: %G?' -1
\end{terminal}
\begin{saida}
Pessoa Qualquer <qualquer@exemplo.com> | assinatura: N
\end{saida}

O Git aceitou a identidade falsa, e o \arq{N} diz que não há assinatura.

\begin{exercicio}{6.2 — Assinar de verdade}
Siga o capítulo \ref{cap:assinatura} com o método que preferir (GPG ou SSH):
\begin{enumerate}
  \item crie ou reaproveite uma chave;
  \item cadastre a chave no GitHub ou no GitLab;
  \item configure o Git para assinar todo commit;
  \item faça um commit num repositório seu e envie com \cmd{git push};
  \item confira o selo na página de commits.
\end{enumerate}
\end{exercicio}

\begin{verifique}
No terminal, \cmd{git log --format='\%h \%G? \%s' -1} mostra a letra \arq{G}. No site, o commit tem o selo \textbf{Verified}. Se aparecer \textbf{Unverified}, confira primeiro o e-mail: \cmd{git config user.email} precisa ser um e-mail verificado na sua conta.
\end{verifique}

\begin{exercicio}{6.3 — Uma tag assinada}
Crie e confira uma tag assinada:
\end{exercicio}

\begin{terminal}
git tag -s v1.0 -m "Versão 1.0"
git tag -v v1.0
git push origin v1.0
\end{terminal}

\begin{verifique}
O \cmd{git tag -v} mostra a mesma confirmação de assinatura do \cmd{git log --show-signature}, e a tag aparece como \textbf{Verified} na página de tags do GitHub.
\end{verifique}

\chapter{Para levar daqui}
\label{cap:memorizar}

\section{Os dois problemas, lado a lado}

Os dois problemas parecem diferentes, mas se completam:

\begin{center}
\begin{tikzpicture}[p/.style={caixa, minimum width=3.6cm, font=\small}]
  \node[p, fill=azulMB, text=white, font=\small\bfseries] (s) at (0,0) {SISTEMA};
  \node[p=verdeMB] (v) at (-3.8,-1.3) {\textbf{Versionamento}\\\scriptsize Problema 1};
  \node[p=ouroMB!85!black] (k) at (3.8,-1.3) {\textbf{Conhecimento}\\\scriptsize Problema 2};
  \node[p=verdeMB] (g) at (-3.8,-2.6) {Git + GitLab};
  \node[p=ouroMB!85!black] (i) at (3.8,-2.6) {IA + validação humana};
  \node[p=verdeMB] (gg) at (-3.8,-3.8) {Histórico e revisão};
  \node[p=ouroMB!85!black] (d) at (3.8,-3.8) {Documentação};
  \node[p, fill=verdeClaro, minimum width=5.4cm] (f) at (0,-5.1) {\textbf{Manutenção e evolução seguras}};
  \draw[setafina] (s) -| (v); \draw[setafina] (s) -| (k);
  \draw[setafina] (v)--(g); \draw[setafina] (g)--(gg); \draw[setafina] (k)--(i); \draw[setafina] (i)--(d);
  \draw[setafina] (gg) |- (f); \draw[setafina] (d) |- (f);
\end{tikzpicture}
\end{center}

\section{Conceitos fundamentais}

\begin{center}
\renewcommand{\arraystretch}{1.45}
\begin{tabularx}{\linewidth}{@{}>{\bfseries\color{azulMB}}l X@{}}
\toprule
Git & Sistema de controle de versão. \\
GitLab & Plataforma de colaboração e gerenciamento baseada em Git. \\
Branch & Linha de desenvolvimento. \\
Commit & Registro de uma alteração no histórico. \\
Merge & Integração de históricos ou alterações. \\
Merge Request & Solicitação para integrar uma branch em outra, normalmente com revisão. \\
GitLab CLI & Ferramenta para interagir com o GitLab pelo terminal (\texttt{glab}). \\
Conflito & Situação em que o Git não consegue decidir sozinho como combinar alterações. \\
IA & Ferramenta auxiliar de análise, desenvolvimento e documentação, sob orientação e validação humana. \\
\bottomrule
\end{tabularx}
\end{center}

\section{Fluxos recomendados}

\begin{center}
\begin{tikzpicture}[p/.style={caixa, minimum width=4.2cm, minimum height=0.62cm, font=\footnotesize, inner sep=2pt}]
  \node[font=\small\bfseries\color{verdeMB}] at (0,0.8) {Para um novo desenvolvimento};
  \foreach \t [count=\i] in {Atualizar o projeto, Criar uma branch, Desenvolver, Testar, Commit, Push, Merge Request, Revisão, Aprovação, Merge}{
    \node[p=verdeMB] (a\i) at (0,-\i*0.82+0.82) {\i.\; \t};}
  \foreach \i [evaluate=\i as \j using int(\i+1)] in {1,...,9}{\draw[setafina] (a\i) -- (a\j);}
  \node[font=\small\bfseries\color{ouroMB!80!black}] at (7,0.8) {Para um sistema legado};
  \foreach \t [count=\i] in {Descobrir, Versionar, Analisar, Documentar, Validar, Corrigir ou refatorar, Testar, Manter}{
    \node[p=ouroMB!85!black] (b\i) at (7,-\i*0.82+0.82) {\i.\; \t};}
  \foreach \i [evaluate=\i as \j using int(\i+1)] in {1,...,7}{\draw[setafina] (b\i) -- (b\j);}
\end{tikzpicture}
\end{center}

\section{Mensagem final}

O conhecimento técnico de uma equipe não deve permanecer apenas na cabeça de quem desenvolveu o sistema.

\begin{itemize}
  \item O código deve ter \textbf{histórico}.
  \item O sistema deve ter \textbf{documentação}.
  \item As alterações devem ter \textbf{revisão}.
  \item As decisões importantes devem ter \textbf{contexto}.
\end{itemize}

As ferramentas de IA podem acelerar esse trabalho, desde que usadas dentro de um processo controlado, com objetivos claros, restrições bem definidas e validação humana.

\begin{center}
\begin{tikzpicture}[p/.style={caixa, minimum width=2.5cm, text width=2.3cm, minimum height=1.2cm, font=\footnotesize}]
  \foreach \t [count=\i] in {Código, Histórico, Documentação, Processo, {Conhecimento compartilhado}}{
    \node[p] at (\i*3.1,0) {\t};}
  \foreach \i in {1,...,4}{\node[font=\large\bfseries\color{ouroMB}] at (\i*3.1+1.55,0) {+};}
  \node[caixa=verdeMB, minimum width=6cm, font=\small\bfseries] at (9.3,-1.5) {= Sistema mais sustentável};
\end{tikzpicture}
\end{center}

\vspace{4mm}
\begin{center}
\begin{minipage}{\linewidth}\centering
{\titulofonte\fontsize{24}{28}\selectfont\color{azulMB} TORNAR OS SISTEMAS MAIS FÁCEIS DE COMPREENDER,\\ MANTER, COMPARTILHAR E EVOLUIR.}
\end{minipage}
\end{center}

\appendix
\chapter{Cola de comandos}
\label{ap:cola}

\renewcommand{\arraystretch}{1.2}

\small
\subsection*{\color{verdeMB}Configuração}
\begin{tabularx}{\linewidth}{@{}>{\leavevmode\ttfamily\small\color{azulMB}}p{6.4cm} X@{}}
git config --global user.name "..." & seu nome \\
git config --global user.email "..." & seu e-mail \\
git config --list & ver a configuração \\
\end{tabularx}

\subsection*{\color{verdeMB}Começar}
\begin{tabularx}{\linewidth}{@{}>{\leavevmode\ttfamily\small\color{azulMB}}p{6.4cm} X@{}}
git init -b main & criar repositório \\
git clone <URL> & copiar um remoto \\
\end{tabularx}

\subsection*{\color{verdeMB}Registrar}
\begin{tabularx}{\linewidth}{@{}>{\leavevmode\ttfamily\small\color{azulMB}}p{6.4cm} X@{}}
git status & estado atual \\
git diff & o que mudou \\
git add arquivo & preparar \\
git commit -m "..." & registrar \\
git commit -am "..." & preparar e registrar os já acompanhados \\
git log --oneline --graph --all & histórico desenhado \\
\end{tabularx}

\subsection*{\color{verdeMB}Branches}
\begin{tabularx}{\linewidth}{@{}>{\leavevmode\ttfamily\small\color{azulMB}}p{6.4cm} X@{}}
git branch & listar \\
git switch -c nome & criar e entrar \\
git switch nome & trocar \\
git branch -m antigo novo & renomear \\
git merge nome & trazer \texttt{nome} para a atual \\
git merge --abort & cancelar o merge \\
\end{tabularx}

\subsection*{\color{verdeMB}Conflitos}
\begin{tabularx}{\linewidth}{@{}>{\leavevmode\ttfamily\small\color{azulMB}}p{6.4cm} X@{}}
git restore --ours arq & manter o lado da branch atual \\
git restore --theirs arq & manter o lado que chega \\
git add arq & marcar como resolvido \\
\end{tabularx}

\subsection*{\color{verdeMB}Remoto}
\begin{tabularx}{\linewidth}{@{}>{\leavevmode\ttfamily\small\color{azulMB}}p{6.4cm} X@{}}
git remote -v & listar remotos \\
git remote add origin <URL> & adicionar remoto \\
git fetch origin & baixar sem mesclar \\
git pull & baixar e mesclar \\
git push origin branch & enviar \\
\end{tabularx}

\subsection*{\color{verdeMB}Desfazer}
\begin{tabularx}{\linewidth}{@{}>{\leavevmode\ttfamily\small\color{azulMB}}p{6.4cm} X@{}}
git restore arq & descartar mudança não preparada \\
git restore --staged arq & tirar da preparação \\
git revert <commit> & desfazer com novo commit \\
\end{tabularx}

\subsection*{\color{verdeMB}Assinatura}
\begin{tabularx}{\linewidth}{@{}>{\leavevmode\ttfamily\small\color{azulMB}}p{6.4cm} X@{}}
gpg --full-generate-key & criar chave GPG \\
gpg --list-secret-keys --keyid-format=long & ver as chaves GPG \\
gpg --armor --export ID & exportar a chave pública \\
git config --global gpg.format ssh & assinar com SSH \\
git config --global user.signingkey ... & escolher a chave \\
git config --global commit.gpgsign true & assinar todo commit \\
git commit -S -m "..." & assinar um commit \\
git log --show-signature -1 & conferir a assinatura \\
git log --format='\%h \%G? \%s' & estado da assinatura (G, N, B, E) \\
git tag -s v1.0 -m "..." & criar tag assinada \\
\end{tabularx}

\subsection*{\color{verdeMB}GitLab CLI}
\begin{tabularx}{\linewidth}{@{}>{\leavevmode\ttfamily\small\color{azulMB}}p{6.4cm} X@{}}
glab auth login & autenticar \\
glab mr create & criar MR \\
glab mr list & listar MRs \\
glab mr view & ver o MR \\
glab mr merge & fazer o merge \\
glab ci status & ver o pipeline \\
\end{tabularx}


% -----------------------------------------------------------------------------
\chapter{Glossário}
\label{ap:glossario}

\begin{description}[style=nextline, font=\bfseries\color{azulMB}, leftmargin=0pt, itemsep=4pt]
\item[Agente de IA] Programa baseado em modelo de linguagem que executa tarefas num ambiente (ler arquivos, rodar comandos, escrever documentos), e não só conversa.
\item[Ancestral comum] Commit a partir do qual duas branches se separaram. Sem ele, o Git recusa o merge comum.
\item[Área de preparação (\emph{staging})] Onde ficam as mudanças escolhidas para o próximo commit.
\item[Assinatura de commit] Marca criptográfica, feita com a sua chave privada, que prova quem criou o commit. Gera o selo \emph{Verified}.
\item[Bare repository] Repositório sem pasta de trabalho, só com o histórico. É o que um servidor guarda.
\item[Branch] Linha de desenvolvimento; tecnicamente, um ponteiro para um commit.
\item[Branch protegida] Branch que não aceita push direto; mudanças só por Merge Request.
\item[CI/CD] Integração e entrega contínuas: testes e implantação automáticos a cada mudança.
\item[Chave pública e chave privada] Par de chaves: a privada fica só com você e assina; a pública pode ser distribuída e serve para conferir a assinatura.
\item[Commit] Registro de um estado do projeto, com autor, data e mensagem.
\item[Conventional Commits] Especificação de mensagens de commit no formato \texttt{tipo(escopo): descrição}.
\item[Conflito] Quando duas mudanças na mesma região de um arquivo não podem ser combinadas automaticamente.
\item[Fast-forward] Merge em que o Git só avança o ponteiro, porque não houve trabalho paralelo.
\item[Gitmoji] Convenção que associa um emoji a cada tipo de mudança, escrito como código (\texttt{:sparkles:}).
\item[GPG] GnuPG, programa que cria chaves e assina dados. Uma das duas formas de assinar commits.
\item[GAP] Lacuna entre o estado atual (AS-IS) e o desejado (TO-BE), ou informação que falta e não deve ser inventada.
\item[HEAD] Referência para onde você está agora (em geral, a branch atual).
\item[Merge] Integração de duas linhas de histórico.
\item[Merge Request (MR)] Pedido para integrar uma branch em outra, com revisão. No GitHub: \emph{Pull Request}.
\item[\texttt{origin}] Nome padrão do repositório remoto principal.
\item[\texttt{ours} / \texttt{theirs}] Num merge: o lado da branch atual / o lado que está chegando.
\item[Pipeline] Sequência automática de etapas (compilar, testar, publicar) executada pelo GitLab.
\item[Prompt] Instrução dada a um modelo ou agente de IA.
\item[Remoto] Cópia do repositório em outro lugar, como o GitLab.
\item[Repositório] Pasta cujo histórico é controlado pelo Git.
\item[Self-Managed] GitLab instalado e administrado pela própria organização.
\item[Token] Na IA, pedaço de texto usado para medir custo e limite. No GitLab, credencial de acesso (trate como senha).
\item[Verified / Unverified] Selos do GitHub e do GitLab: assinatura conferida, ou assinatura presente mas que não pôde ser conferida.
\end{description}

% -----------------------------------------------------------------------------
\chapter{Plano da apresentação}
\label{ap:plano}

Roteiro de tempo da capacitação de 01/10/2026, das 10h às 11h. Os números de slide se referem a \arq{apresentacao/apresentacao.pdf}.

\begin{center}
\renewcommand{\arraystretch}{1.4}
\begin{tabularx}{\linewidth}{@{}>{\bfseries\color{azulMB}}l c c >{\raggedright\arraybackslash}X@{}}
\toprule
\textbf{Horário} & \textbf{Duração} & \textbf{Slides} & \textbf{Conteúdo} \\
\midrule
10h00 & 3 min & 1--2 & Abertura, agenda e objetivos \\
10h03 & 19 min & 3--15 & Fundamentos: por que versionar, Git, GitLab, Self-Managed, branch, Merge Request, comandos, fluxo, conflito, \texttt{glab} \\
10h22 & 15 min & 16--25 & Problema 1: cenário, históricos sem relação, os oito passos, \texttt{ours} e \texttt{theirs} \\
10h37 & 18 min & 26--35 & Problema 2: IA como ferramenta, evolução dos prompts, falhas e regras, números, fases, travas, idioma, entregas \\
10h55 & 5 min & 36--39 & Encerramento, mensagem final e perguntas \\
\bottomrule
\end{tabularx}
\end{center}

\begin{dica}
Se o tempo apertar, os slides 8 (recursos do GitLab), 12 (tabela de comandos) e 35 (o que o prompt entrega) podem ser passados rapidamente: o conteúdo deles está completo nesta apostila.
\end{dica}

% -----------------------------------------------------------------------------
\chapter{Créditos e licenças}
\label{ap:creditos}

Todo o material de terceiros usado aqui está salvo localmente, na pasta \arq{assets/} (imagens e fontes), e o manual da marca em \arq{materials/manual-da-marca/}, para que a apostila e a apresentação possam ser recompiladas sem internet.

\begin{center}
\renewcommand{\arraystretch}{1.35}
\footnotesize
\begin{tabularx}{\linewidth}{@{}>{\raggedright\arraybackslash}p{3.4cm} >{\raggedright\arraybackslash}X >{\raggedright\arraybackslash}p{3.3cm}@{}}
\toprule
\textbf{Item} & \textbf{Origem} & \textbf{Licença} \\
\midrule
Manual de Identidade Visual da Marinha do Brasil (cores e tipografia) & CCEM, dez/2025 (\arq{materials/manual-da-marca/}) & Material institucional da Marinha do Brasil \\
Foto da Ilha das Cobras & Helder da Rocha, via Wikimedia Commons (\enquote{IlhaCobras.jpg}), reduzida & CC BY-SA 2.0 \\
Logotipo do Git & Jason Long, via Wikimedia Commons & CC BY 3.0 \\
Logotipo do GitLab (tanuki) & GitLab Inc., \arq{gitlab-com/gitlab-artwork} & CC BY-NC-SA 4.0; marca registrada da GitLab Inc. \\
Símbolo do GitHub & GitHub Octicons, via Wikimedia Commons & MIT; marca da GitHub Inc. \\
Símbolo do Claude & Wikimedia Commons & CC0; marca da Anthropic \\
Tirinha xkcd \#1597 \enquote{Git} & Randall Munroe, \url{https://xkcd.com/1597/} & CC BY-NC 2.5 \\
Fonte Montserrat & Julieta Ulanovsky e colaboradores & SIL Open Font License 1.1 \\
Fonte Bebas Neue & Dharma Type, via Google Fonts & SIL Open Font License 1.1 \\
Fonte Fira Mono & Mozilla / Carrois & SIL Open Font License 1.1 \\
\bottomrule
\end{tabularx}
\end{center}

As marcas de terceiros (Git, GitLab, GitHub, Claude) aparecem apenas para identificar as ferramentas de que o texto trata. As licenças NC (não comercial) são compatíveis com o uso deste material: capacitação interna, sem fins lucrativos.

\subsection*{Identidade visual}

As cores seguem o Manual de Identidade Visual da Marinha do Brasil: Azul Marinha \texttt{\#050F41}, Amarelo Ouro \texttt{\#FAB932} e Verde Brasil \texttt{\#079551}. A tipografia de apoio oficial é a Montserrat, usada no texto. As tipografias de título do manual (Octin College e Bebas Neue Pro) são comerciais. No lugar delas, usamos a Bebas Neue livre, do mesmo desenho da alternativa oficial.

% -----------------------------------------------------------------------------
\chapter{Para continuar estudando}
\label{ap:referencias}

\begin{itemize}
  \item \textbf{Pro Git}, de Scott Chacon e Ben Straub. O livro de referência do Git, gratuito e com tradução para o português: \url{https://git-scm.com/book/pt-br/v2}.
  \item \textbf{Documentação oficial do Git}: \url{https://git-scm.com/docs}.
  \item \textbf{Documentação do GitLab}: \url{https://docs.gitlab.com}. Veja em especial \emph{Merge requests} e \emph{Protected branches}.
  \item \textbf{GitLab CLI (\texttt{glab})}: \url{https://gitlab.com/gitlab-org/cli}.
  \item \textbf{Claude Code}: \url{https://docs.claude.com/en/docs/claude-code/overview}.
  \item \textbf{Conventional Commits}: \url{https://www.conventionalcommits.org/pt-br/v1.0.0/}.
  \item \textbf{Gitmoji}: \url{https://gitmoji.dev}.
  \item \textbf{Padrões de commits}, de Iuri Silva, uma referência em português: \url{https://github.com/iuricode/padroes-de-commits}.
  \item \textbf{Assinatura de commits no GitHub}: \url{https://docs.github.com/pt/authentication/managing-commit-signature-verification}.
  \item \textbf{Assinatura de commits no GitLab}: \url{https://docs.gitlab.com/user/project/repository/signed_commits/}.
  \item \textbf{Learn Git Branching}, exercícios visuais e interativos de branch e merge, em português: \url{https://learngitbranching.js.org/?locale=pt_BR}.
\end{itemize}

\subsection*{Arquivos deste material}

\begin{center}
\renewcommand{\arraystretch}{1.3}
\small
\begin{tabularx}{\linewidth}{@{}>{\ttfamily\footnotesize}l X@{}}
\toprule
apresentacao/apresentacao.pdf & os slides da capacitação \\
apostila/apostila.pdf & esta apostila \\
model/SYSTEM\_DOCUMENTATION.md & o prompt final (V4) de documentação \\
roteiro\_versao\_01.md & o roteiro inicial, com os assuntos obrigatórios \\
roteiro\_versão\_final.md & o roteiro final \\
\bottomrule
\end{tabularx}
\end{center}

\end{document}
